\documentclass[11pt]{article}
\usepackage[a4paper,margin=27mm]{geometry}
\usepackage[T1]{fontenc}
\usepackage{lmodern,microtype,amsmath,amssymb,amsthm,mathtools,enumitem}
\usepackage[hidelinks,hypertexnames=false]{hyperref}
\newtheorem{theorem}{Theorem}
\newtheorem{lemma}[theorem]{Lemma}
\newtheorem{proposition}[theorem]{Proposition}
\newtheorem{corollary}[theorem]{Corollary}
\theoremstyle{definition}\newtheorem{definition}[theorem]{Definition}
\theoremstyle{remark}\newtheorem{remark}[theorem]{Remark}
\newcommand{\R}{\mathbb R}
\newcommand{\Dscr}{\mathcal D}
\newcommand{\supp}{\operatorname{supp}}
\newcommand{\debt}{\mathrm{debt}}
\newcommand{\cred}{\mathrm{credit}}
\newcommand{\sil}{\mathrm{silence}}
\newcommand{\Band}{\mathrm{B}}
\newcommand{\Jid}{\mathcal J}
\newcommand{\Kid}{\mathcal K}
\newcommand{\Uf}{\mathcal U}
\newcommand{\Lo}[3]{\mathrm L^{#1}_{#2}(#3)}
\newcommand{\Hi}[3]{\mathrm H^{#1}_{#2}(#3)}
\newcommand{\CN}{\mathrm{CN}}
\newcommand{\DI}{\mathrm{DI}}
\title{Fineness Spreads Along Edges}
\author{}\date{11 October 2026}
\begin{document}\maketitle

\noindent
This note continues CCF (\emph{Closed Choice Forces Fineness}).
It also cites CBOI (\emph{Closed Choice Beyond Ordinary Integrals}), EDL
(\emph{Edge Dichotomy and Leverage}), FAED (\emph{Finitely Additive Edge
Dichotomy}) and IC (\emph{Independence Chains and Two Additive Worlds}). It
concerns
the conjecture for complete rankings, problem~P1 of
Section~\ref{sec:problem}: three worlds, every ranking a complete preorder
satisfying (G) and (IB), and (S). By CCF, a world whose debt and credit
orders are \emph{fine} has a finitely additive weak extension of CBOI's
form. We show that fineness at one world spreads to all worlds, that
fineness on one edge already makes the vacuous model $V$ dominated, and we
collect what a counterexample must look like.
\begin{enumerate}[leftmargin=*,itemsep=3pt]
\item \emph{Edge transfer} (Theorem~\ref{thm:transfer},
Corollary~\ref{cor:edgefine}). Let $n\ge3$, assume (G) at every world and
(S), and let worlds $1$ and $j$ be complete with (IB). If world $j$ is fine
on the region $D_{1j}$ where worlds $1$ and $j$ disagree, then so is
world~$1$. More generally, if world~$j$ has a finitely additive bound of
CBOI's form on $D_{1j}$, then every non-fine part of world~$1$ on $D_{1j}$
sits at an atom of world~$j$'s charges. Nothing is assumed about~$V$.
\item \emph{One fine world} (Theorem~\ref{thm:onefine}). Let $n=3$, let all
rankings be complete with (G) and (IB), and assume (S). If one world is fine,
all three are, and $V$ is dominated.
\item \emph{One fine edge} (Theorem~\ref{thm:fineedge},
Corollary~\ref{cor:fineedge}). In the same setting, if two worlds are fine
on the region where they disagree, then $V$ is dominated.
\item \emph{One tame world} (Theorem~\ref{thm:tame},
Theorem~\ref{thm:aext}). In the same setting, $V$ is dominated as soon as
one world is fine, or has a strictly positive one-sided loss of CBOI's form,
or a nonconstant additive primary with values in an ordered group, or
satisfies all finite cancellation conditions. The same holds if one
world has strictly positive two-sided finitely additive extensions on both of
its edge regions after removing sets of arbitrarily small measure.
\item \emph{Low-cone heaviness} (Theorem~\ref{thm:LH}). In the same
setting, with no assumption on $V$: at every world, credit on any non-null
set compensates some non-null debt, and debt on some part of any non-null set
is not compensated by some non-null credit.
\item \emph{Archimedean axioms} (Section~\ref{sec:arch}). Germ continuity at
one world, which conic continuity implies, gives fineness, and so $V$ is
dominated. An example shows that fineness does not follow from completeness,
(G) and (IB), even for a world ranked exactly by a finitely additive loss.
\item \emph{Shape of a counterexample} (Theorem~\ref{thm:shape}). No world
is fine. No world carries a strictly positive one-sided finitely additive
loss, so every world carries a leverage sequence and violates a finite
cancellation condition. No edge is fine, and every edge carries a selection
that omits both endpoints. (LH) and its dual hold at every world. At every
such selection both edge worlds are densely credit-heavy on the selection's
high cones, and their debt orders are non-fine on their own odd regions
relative to credit on deep tails of the third region.
\end{enumerate}

\emph{What remains open.} Problem~P1 itself remains open: we do not know
whether a counterexample exists. Section~\ref{sec:shape} states the residual
precisely. P1 would be settled, with $V$ dominated, by a positive answer
to the first of the following questions or by a negative answer to the
second: whether (S) forces, at one world,
the relative fineness (M2$^*$) of Theorem~\ref{thm:ownregion}; and whether a
complete (G)+(IB) order can fail to have finitely additive extensions off
small sets without having a nonconstant group-valued primary
(Remark~\ref{rem:Q}). The case $n\ge4$ of Theorem~\ref{thm:onefine} and the
case of incomplete rankings are also open.

The results are taken from working notes that were checked by independent
referees in four parts: the edge transfer, one fine world and the
Archimedean axioms (Sections~\ref{sec:transfer} and~\ref{sec:arch}), the
fine edge (Section~\ref{sec:edge}), the tame world
(Section~\ref{sec:tame}), and low-cone heaviness (Section~\ref{sec:LH}).
Their corrections are incorporated. Results that the
working notes took from earlier, unpublished working notes are stated and
proved in full here.

\tableofcontents

\section{Setting}\label{sec:setting}

\subsection{Problems, models and axioms}

We use the setting of CCF, which is that of CBOI, EDL, FAED and IC.
There are $n\ge3$ worlds, $\Delta\coloneqq\{p\in\R^n:p\ge0,\ \sum_ip_i=1\}$,
and $e_i$ is the $i$-th vertex. Binary problems $\{g,-g\}$ are represented
by $g\in H\coloneqq\{g\in\R^n:g_1>0\}$, with normalized Gaussian
measure~$\mu$; coordinate hyperplanes are null. Models are measurable and
identified modulo null sets. At world~$i$ a model's advice at $g$ has a
guidance rank $r_i(g,M)\in\{0,1,2,3\}$ (good, silent, bad, contradictory).
A model is \emph{regular} if it is never contradictory and correct on
unanimous problems. Coherent models $\Dscr_K$ ($K\subseteq\Delta$ nonempty,
compact, convex) are regular, $V\coloneqq\Dscr_\Delta$, and $D_i$ denotes
the ideal model of world~$i$, which recommends the option good at~$i$ on
every problem. $H'$ is the set of mixed problems in $H$.

At each world $\preceq_i$ is a preorder on models; smaller is better.
$X\prec_iY$ means $X\preceq_iY$ and not $Y\preceq_iX$; so $X\prec_iY$
excludes $Y\preceq_iX$ without any completeness. (G) is strict guidance
monotonicity and (IB) is independence under replacement of shared advice.
(S) supplies coherent undominated models $C_p$, $p\in\Delta$, competitive on
$\supp p$, with closed sections $A_i(X)\coloneqq\{p:C_p\preceq_iX\}$. A model
$C$ is \emph{competitive} on $S$ if for every model $X$ some $i\in S$ has
$C\preceq_iX$. A model $M$ is \emph{dominated} if some model $X$ has
$X\preceq_kM$ at every world and $X\prec_kM$ at some world. By CBOI's
Lemma~1 every model has a selected coherent weak dominator at all worlds, so
a dominated model is dominated by a selected coherent model. If $V$ is
undominated, it is competitive on all worlds (CCF, Section~1).

By the face lemma (recalled in CCF, Section~1; for edges see
FAED, Section~1), the credal set of $C_p$ lies in the face spanned by
$\supp p$; in particular $C_{e_i}=D_i$. We call $C_p$ a \emph{selection}, and
a selection \emph{on the edge} $[e_v,e_w]$ if $p\in[e_v,e_w]$; if $p$ lies in
the open edge $(e_v,e_w)$, $C_p$ is competitive on $\{v,w\}$.

We use two consequences of (S) (CCF, Section~3).
\begin{enumerate}[label=(T\arabic*),leftmargin=*,itemsep=2pt]
\item If world $k$ is complete, $Y$ is regular and $Y\prec_kC_{p^*}$, then
$Y\prec_kC_p$ for all $p$ near $p^*$. [$p^*\notin A_k(Y)$, which is closed;
near $p^*$, $C_p\not\preceq_kY$, and completeness gives $Y\prec_kC_p$.]
\item If $p^*\in[e_v,e_w]$, $F$ is regular and $F\prec_wC_{p^*}$, then
$C_p\preceq_vF$ for all $p$ in the open edge $(e_v,e_w)$ near $p^*$.
[$C_p\not\preceq_wF$ near $p^*$, and $C_p$ is competitive on $\{v,w\}$.]
\end{enumerate}

\subsection{Debt, credit and pairs}\label{sec:pairs}

On a mixed problem a regular model gives the option good at world $v$
(\emph{credit} at $v$, rank $0$), silence (rank $1$), or the option bad
at~$v$ (\emph{debt} at $v$, rank $2$). For a regular $M$ write
$\cred_v(M)$, $\debt_v(M)$ and $\sil(M)$ for the corresponding subsets of
$H'$; for $v=1$ the subscript is dropped, as in CCF. For measurable
$Q\subseteq H'$, $M+\debt_v(Q)$ and $M+\cred_v(Q)$ agree with $M$ except that
they give debt, resp.\ credit, at $v$ on $Q$, and ``$M$ silenced on $Q$''
gives silence there. All these models are regular.

For measurable $B,C\subseteq H'$, following CCF, Section~3,
\begin{gather*}
 B\preceq^v_dC\ :\Longleftrightarrow\ V+\debt_v(B)\preceq_vV+\debt_v(C),\\
 B\preceq^v_cC\ :\Longleftrightarrow\ V+\cred_v(C)\preceq_vV+\cred_v(B).
\end{gather*}
So $B\preceq^v_dC$ says that debt on $B$ is no worse than debt on $C$, and
$B\preceq^v_cC$ says that credit on $B$ is worth no more than credit on
$C$. Both are transitive, and complete if $\preceq_v$ is. By (IB), $V$ may
be replaced by any regular model silent on $B\cup C$, and the comparison
depends only on $B\setminus C$ and $C\setminus B$; overlapping sets may be
compared. By (G), $B\subseteq C$ implies $B\preceq^v_dC$ and
$B\preceq^v_cC$; if moreover $C\setminus B$ is non-null, both are strict.
For world~$1$ we write $\preceq_d,\preceq_c$.

A \emph{pair} $(A,T)$ at world $v$ consists of disjoint measurable
$A,T\subseteq H'$ (debt part and credit part). It is \emph{SI} if
$V+\debt_v(A)+\cred_v(T)\prec_vV$, \emph{weakly SI} if
$V+\debt_v(A)+\cred_v(T)\preceq_vV$, \emph{AC} if it is not SI, and
\emph{strictly AC} if $V\prec_vV+\debt_v(A)+\cred_v(T)$. If world $v$ is
complete, AC means $V\preceq_vV+\debt_v(A)+\cred_v(T)$. By (IB), $V$ may be
replaced by any regular model silent on $A\cup T$. Pairwise disjoint SI pairs
compose to an SI pair, and, if world $v$ is complete, pairwise disjoint
AC pairs to an AC pair (CCF, Section~6). By strict (G), an SI pair has non-null credit part, a pair
$(A,\varnothing)$ with $A$ non-null is strictly AC, and a pair
$(\varnothing,T)$ with $T$ non-null is SI.

\subsection{Fineness, germ points and ultrafilters}\label{sec:germs}

\begin{definition}\label{def:fine}
Let $A,T\subseteq H'$ be non-null. $\Jid^v_A$ is the family of measurable
$B\subseteq H'$ contained in a finite union of measurable sets each
$\preceq^v_dA$; $\Kid^v_T$ is defined likewise with $\preceq^v_c$. For a
non-null measurable $R\subseteq H'$, world $v$'s debt order is \emph{fine on
$R$} if $R\in\Jid^v_A$ for every non-null $A\subseteq R$, and its credit
order is fine on $R$ if $R\in\Kid^v_T$ for every non-null $T\subseteq R$.
World $v$ is \emph{fine on $R$} if both orders are, and \emph{fine} if it is
fine on $H'$. For $v=1$ we write $\Jid_A,\Kid_T$.
\end{definition}

For $R=H'$ this is CCF's notion: a finite union of sets $\preceq A$ can be
made a partition by (G). Fineness on $R$ uses only yardsticks inside $R$; by
Lemma~\ref{lem:fineregion} below, for an edge region it is the same as
fineness with yardsticks anywhere.

\begin{lemma}\label{lem:ideal}
$\Jid^v_A$ and $\Kid^v_T$ are ideals of the measure algebra: they contain
the null sets and are closed under measurable subsets and finite unions. If
$A'\preceq^v_dA$ then $\Jid^v_{A'}\subseteq\Jid^v_A$; if
$T'\preceq^v_cT$ then $\Kid^v_{T'}\subseteq\Kid^v_T$.
\end{lemma}
\begin{proof}
Null sets are equivalent to $\varnothing$, which is $\preceq A$ by (G). If
$B\subseteq B_1\cup\dots\cup B_m$ with $B_l\preceq^v_dA$, every measurable
subset of $B$ lies in the same union, and a finite union of such sets lies
in the union of the covers. The last statements follow by transitivity. The
credit case is the same.
\end{proof}

For an edge $[e_1,e_j]$ write $D_{1j}\coloneqq\{g\in H':g_j<0\}$ (CCF writes
$D_j$ for this set; we reserve $D_j$ for the ideal model) and
$t_j\coloneqq g_1/|g_j|$ on $D_{1j}$. For $0\le s<s'\le\infty$ put
$\Band_j(s,s')\coloneqq D_{1j}\cap\{s<t_j<s'\}$, and similarly with closed
ends; level sets of $t_j$ are null, so the ends do not matter.

\begin{definition}\label{def:germ}
A \emph{germ point} of the edge $[e_1,e_j]$ is a pair $(z,\sigma)$ with
$z\in[0,\infty]$ and $\sigma\in\{-,+\}$, where $\sigma=+$ if $z=0$ and
$\sigma=-$ if $z=\infty$. Its \emph{one-sided bands} are
$\Band_j(z-\varepsilon,z)$ for $\sigma=-$ (with $\Band_j(1/\varepsilon,\infty)$
if $z=\infty$) and $\Band_j(z,z+\varepsilon)$ for $\sigma=+$, for
$\varepsilon>0$. The germ point is \emph{$A$-heavy} if none of its one-sided
bands lies in $\Jid_A$, and \emph{$T$-heavy} if none lies in $\Kid_T$.
\end{definition}

The bands decrease as $\varepsilon\downarrow0$, so by Lemma~\ref{lem:ideal}
a germ point fails to be $A$-heavy as soon as one of its bands lies in
$\Jid_A$.

\begin{lemma}\label{lem:germ}
If $D_{1j}\notin\Jid_A$, some germ point of $[e_1,e_j]$ is $A$-heavy. If
$D_{1j}\notin\Kid_T$, some germ point is $T$-heavy. No (S) is used.
\end{lemma}
\begin{proof}
Suppose no germ point is $A$-heavy. For $z\in(0,\infty)$ there are
$\varepsilon,\varepsilon'>0$ with $\Band_j(z-\varepsilon,z)$ and
$\Band_j(z,z+\varepsilon')$ in $\Jid_A$, so their union, which is the band
$\Band_j(z-\varepsilon,z+\varepsilon')$ up to a null set, is in $\Jid_A$.
The points $0$ and $\infty$ have neighbourhoods $[0,\varepsilon)$ and
$(1/\varepsilon,\infty]$ whose bands are in $\Jid_A$. Finitely many of these
open sets cover the compact space $[0,\infty]$, so $D_{1j}$ is, up to a null
set, a finite union of members of $\Jid_A$, hence in $\Jid_A$. The credit
case is the same.
\end{proof}

All of this is stated in world~$1$'s frame. For distinct worlds $v,w$ put
\[
 D_{vw}\coloneqq\{g\in H':g_vg_w<0\},\qquad
 t^v_w\coloneqq|g_v|/|g_w|\ \text{on }D_{vw},
\]
so that $D_{1j}$ is as above, $t^1_j=t_j$ and $t^w_v=1/t^v_w$. The law of
$t^v_w$ under $\mu|_{D_{vw}}$ has a positive density on $(0,\infty)$, and
for distinct $j,k\ne v$ the joint law of $(t^v_j,t^v_k)$ on $D_{vj}\cap D_{vk}$
has a positive density on $(0,\infty)^2$ (CCF, Section~1). The axioms are
symmetric in the labels of the worlds, and every object above is a set of
problems, so every definition and lemma stated for world~$1$ and the edge
$[e_1,e_j]$ holds for world $v$ and the edge $[e_v,e_w]$, with $D_{vw}$ and
$t^v_w$ in place of $D_{1j}$ and $t_j$; the change of representing
half-space is harmless, as in the proof of CCF's Theorem~16. We use this
relabelling throughout without further comment.

Let $\mathfrak A$ be the measure algebra of $H'$: measurable subsets of $H'$
modulo null sets. Ultrafilters below are ultrafilters of $\mathfrak A$. An
ultrafilter $\Uf$ \emph{avoids} an ideal $\Jid$ if $\Uf\cap\Jid=\varnothing$.

\begin{lemma}\label{lem:uf}
\begin{enumerate}[label=(\alph*)]
\item Every ultrafilter contains sets of arbitrarily small measure. If a
finite union of measurable sets is in $\Uf$, one of them is.
\item If $R\notin\Jid_A$, some ultrafilter contains $R$ and avoids $\Jid_A$.
If $(z,\sigma)$ is an $A$-heavy germ point of $[e_1,e_j]$, some ultrafilter
avoids $\Jid_A$ and contains all one-sided bands of $(z,\sigma)$. The same
holds for $\Kid_T$.
\item If $\Uf$ avoids $\Jid_A$ and world~$1$ is complete, then
$A\prec_dB$ for every $B\in\Uf$; and $A'\prec_dB$ whenever moreover
$A'\preceq_dA$. The same holds for $\Kid_T$ and $\prec_c$.
\item If $\Uf$ contains $D_{1j}$, there is a unique germ point of $[e_1,e_j]$
all of whose one-sided bands are in $\Uf$. If $\Uf$ avoids $\Jid_A$ (resp.\
$\Kid_T$), this germ point is $A$-heavy (resp.\ $T$-heavy).
\end{enumerate}
\end{lemma}
\begin{proof}
(a) Split $B\in\Uf$ into two halves of equal measure ($\mu$ is atomless); one
half is in $\Uf$; iterate. The second statement holds in every ultrafilter
of a Boolean algebra.

(b) The family $\{R\}\cup\{H'\setminus J:J\in\Jid_A\}$ has the finite
intersection property: $R\setminus(J_1\cup\dots\cup J_s)$ is non-null, since
otherwise $R\subseteq J_1\cup\dots\cup J_s\in\Jid_A$. It extends to an
ultrafilter, which avoids $\Jid_A$ because it contains $H'\setminus J$ for
each $J\in\Jid_A$. For the germ point use the family of all its bands
together with $\{H'\setminus J\}$: a finite intersection is a band minus a
member of $\Jid_A$, and it is non-null because no band is in $\Jid_A$.

(c) If $B\preceq_dA$ then $B\in\Jid_A$. So $B\not\preceq_dA$, and
completeness gives $A\prec_dB$; then $A'\preceq_dA\prec_dB$.

(d) For $z\in[0,\infty]$ and $\varepsilon>0$ let $N(z,\varepsilon)$ be the
band of $t_j$-levels within $\varepsilon$ of $z$ (one-sided at $0$ and at
$\infty$, with $(1/\varepsilon,\infty]$ as neighbourhood of $\infty$). If
for every $z$ some $N(z,\varepsilon_z)$ were not in $\Uf$, a finite subcover
of $[0,\infty]$ would write $D_{1j}$ as a finite union of non-members, which
is impossible by (a). So some $z$ has $N(z,\varepsilon)\in\Uf$ for all
$\varepsilon$. For $z\in(0,\infty)$, $N(z,\varepsilon)$ is the union of the
two one-sided bands of size $\varepsilon$ up to a null set, so one of them is
in $\Uf$. If the left band of some size is not in $\Uf$, neither is any
smaller left band, so all small right bands, and hence all right bands, are
in $\Uf$. So all bands of one side are in $\Uf$. Uniqueness: the bands of
distinct germ points, including $(z,-)$ and $(z,+)$, are eventually
disjoint. Heaviness: members of $\Uf$ are not in $\Jid_A$.
\end{proof}

\subsection{Three worlds: regions and options}\label{sec:regions}

For $n=3$ let $P_x$ be the set of mixed problems at which world $x$ is the
odd one out:
\[
 P_1=\{g_2<0,\ g_3<0\},\qquad P_2=\{g_2<0<g_3\},\qquad P_3=\{g_3<0<g_2\}
\]
inside $H$. They partition $H'$ up to null sets, and $D_{vw}=P_v\cup P_w$.
As in IC, $o_K$ denotes the option good exactly at the worlds in $K$. On
$P_x$, with $\{x,y,z\}=\{1,2,3\}$, the option $o_{yz}$ is debt at $x$ and
credit at $y$ and $z$, and $o_x$ is credit at $x$ and debt at $y$ and $z$:
\begin{center}
\begin{tabular}{c|ccc|ccc}
 & \multicolumn{3}{c|}{accept $g$} & \multicolumn{3}{c}{accept $-g$}\\
 & option & credit at & debt at & option & credit at & debt at\\\hline
$P_1$ & $o_1$ & $1$ & $2,3$ & $o_{23}$ & $2,3$ & $1$\\
$P_2$ & $o_{13}$ & $1,3$ & $2$ & $o_2$ & $2$ & $1,3$\\
$P_3$ & $o_{12}$ & $1,2$ & $3$ & $o_3$ & $3$ & $1,2$
\end{tabular}
\end{center}
On $D_{vw}$ the ranks of worlds $v$ and $w$ add up to $2$ for every regular
advice, so
\begin{equation}
 \cred_v=\debt_w\quad\text{and}\quad\debt_v=\cred_w\qquad\text{on }D_{vw}.
 \label{eq:swap}
\end{equation}
On each $P_x$ the joint law of two ratios $|g_a|/|g_b|$ and $|g_c|/|g_d|$
with $\{a,b\}\ne\{c,d\}$ has a positive density on $(0,\infty)^2$, since two
such ratios determine the problem up to scale. Hence every subset of $P_x$ defined by a nonempty open
condition on two such ratios is non-null; we use this without comment.

For $w\ne v$ and $0<a,b<\infty$ put
\[
 \Lo vwa\coloneqq D_{vw}\cap\{t^v_w<a\},\qquad
 \Hi vwb\coloneqq D_{vw}\cap\{t^v_w>b\},
\]
the \emph{low cone} and the \emph{high tail} of world $v$ on the edge
$[e_v,e_w]$; $\Lo vwa=\Hi wv{1/a}$.

\subsection{Edge models}\label{sec:edgemodels}

By the face lemma a selection on $[e_v,e_w]$ has a credal interval
$\{se_v+(1-s)e_w:s\in[a,b]\}$. Put $\alpha\coloneqq(1-a)/a\ge
\beta\coloneqq(1-b)/b$, allowing $0$ and $\infty$. Then (EDL, Section~1) the
selection gives credit at $v$ on $D_{vw}\cap\{t^v_w>\alpha\}$, silence on
$D_{vw}\cap\{\beta<t^v_w<\alpha\}$, debt at $v$ on
$D_{vw}\cap\{t^v_w<\beta\}$, and the option good at $v$ and $w$ off
$D_{vw}$. We write it $(\alpha,\beta)$, in world $v$'s frame; in world $w$'s
frame it is $(1/\beta,1/\alpha)$. Regular models of this form, coherent or
not, are \emph{edge models}. The credal set contains $e_v$ iff $\beta=0$ and
$e_w$ iff $\alpha=\infty$. A selection whose credal set contains neither
endpoint lies in the open edge and has $0<\beta\le\alpha<\infty$; we call it
\emph{two-sided}. In world~$1$'s frame on $[e_1,e_j]$ we use CCF's endpoint
families
\[
 E_j(b)\coloneqq(\infty,b),\qquad E^j_1(a)\coloneqq(a,0),
\]
so that $E^j_1(0)=D_1$, $E_j(\infty)=D_j$ and
$E^j_1(\infty)=E_j(0)=V_{1j}\coloneqq\Dscr_{[e_1,e_j]}$. For edge models,
$(\alpha,\beta)\le_{\mathrm{pw}}(\alpha',\beta')$ means $\alpha\le\alpha'$ and
$\beta\le\beta'$; then the first is pointwise at least as good at world $v$,
and pointwise at least as bad at world $w$.

\subsection{Losses of finitely additive form}\label{sec:losses}

Following CBOI and CCF, Section~1, a loss of \emph{CBOI's form} at world $k$
is $L_k(M)\coloneqq\sum_{l=1}^3\nu_{kl}\{g:r_k(g,M)\ge l\}$, where the
$\nu_{kl}$ are finite nonnegative finitely additive measures on the
measurable subsets of $H$ with
\begin{equation}
 \nu_{kl}(A)=0\text{ if }\mu(A)=0,\qquad\nu_{kl}(A)>0\text{ if }\mu(A)>0.
 \label{eq:positive}
\end{equation}
For regular models only $\nu_{k1},\nu_{k2}$ on $H'$ matter. The one-sided
bound of CCF is
\begin{equation}
 B\text{ coherent undominated},\ X\text{ regular},\ B\preceq_kX
 \ \Longrightarrow\ L_k(B)\le L_k(X).
 \tag{W$_k$}\label{eq:W}
\end{equation}
It needs no completeness. By EDL's Theorem~2, world $k$ has a loss of CBOI's
form with \eqref{eq:positive} and \eqref{eq:W} if and only if it carries no
\emph{leverage sequence} in EDL's sense. By CCF's Lemma~1 (CBOI's Theorem~2
in the one-sided form of EDL's Proposition~1), if $n\ge3$, (G) and (S) hold
and every world has such a loss, every coherent model whose credal set
contains three vertices is dominated by a coherent model.

\subsection{The problem}\label{sec:problem}

\begin{quote}
\textbf{P1.} Let $n=3$, let every ranking be a complete preorder satisfying
(G) and (IB), and assume (S). Is $V$ dominated?
\end{quote}
An edge $[e_i,e_j]$ has the \emph{endpoint property} (EP$_{ij}$) if every
selection on it has $e_i$ or $e_j$ in its credal set (CCF, Definition~2). By
CCF's Theorems~7 and~15, a world with (EP) on both of its edges is fine and
has a weak extension of CBOI's form. By CCF's Theorem~18, two worlds with
such extensions, or merely with one-sided bounds \eqref{eq:W}, make $V$
dominated if the third world is complete with (IB). By CCF's Theorem~16, $V$ is
dominated if every edge has (EP).

\section{Tools}\label{sec:tools}

Throughout this section, $j\ge2$, $D\coloneqq D_{1j}$, $t\coloneqq t_j$, and
$C_q=(\alpha_q,\beta_q)$ for $q\in[e_1,e_j]$, in world~$1$'s frame. By the
convention of Section~\ref{sec:germs} every statement also holds for any
world $v$ and edge $[e_v,e_w]$.

\subsection{The vertex blocks of an edge}

\begin{lemma}\label{lem:ends}
Assume (G) at worlds $1$ and $j$ and (S). Let $q^*\in[e_1,e_j]$ and
$q_m\to q^*$ along the edge.
\begin{enumerate}[label=(\alph*)]
\item If $C_{q^*}=D_j$, then $\beta_{q_m}\to\infty$. If $C_{q^*}=D_1$, then
$\alpha_{q_m}\to0$.
\item $S_j\coloneqq\{q\in[e_1,e_j]:C_q=D_j\}$ is closed, contains $e_j$ and
not $e_1$. Let $q_e$ be its point nearest to $e_1$. Every $q\in[e_1,q_e)$
has $C_q\ne D_j$, and there are $q_m\uparrow q_e$ in the open edge.
\item $S_1\coloneqq\{q\in[e_1,e_j]:C_q=D_1\}$ is closed, contains $e_1$ and
not $e_j$. It has a boundary point $q_0$ in the edge, and every boundary
point is a limit of points $q_m$ of the open edge with $C_{q_m}\ne D_1$.
\end{enumerate}
\end{lemma}
\begin{proof}
(a) Suppose $\beta_{q_m}\le b$ along a subsequence, with $1\le b<\infty$, and
put $X\coloneqq E_j(2b)$. Each $C_{q_m}$ satisfies $C_{q_m}\le_{\mathrm{pw}}X$,
so $C_{q_m}\preceq_1X$ by (G), and closedness of $A_1(X)$ gives
$D_j\preceq_1X$. But $X$ is pointwise better than $D_j$ at world~$1$,
strictly on $\Band_j(2b,\infty)$, so $X\prec_1D_j$: a contradiction. If
$\alpha_{q_m}\ge a>0$ along a subsequence, put $X\coloneqq E^j_1(a/2)$. Then
$X\le_{\mathrm{pw}}C_{q_m}$, so $C_{q_m}\preceq_jX$ by (G) (the pointwise order
is reversed at $j$), and closedness of $A_j(X)$ gives $D_1\preceq_jX$. But
$X$ is silent on $\Band_j(0,a/2)$, where $D_1$ gives credit at $1$, which is
debt at $j$; so $X\prec_jD_1$.

(b) For every regular $M$, $D_j\preceq_jM$ by (G), and $M\preceq_jD_j$ only
if $M=D_j$ by strict (G). Hence $S_j=A_j(D_j)\cap[e_1,e_j]$, which is closed.
By the face lemma $C_{e_j}=D_j$ and $C_{e_1}=D_1\ne D_j$. The rest is the
choice of $q_e$. (c) The same with $A_1(D_1)$; a nonempty proper closed subset
of the connected edge has a boundary point $q_0$, and $q_0\ne e_j$, so
nearby points outside $S_1$ lie in the open edge.
\end{proof}

\subsection{Yardsticks}

\begin{lemma}\label{lem:E12}
Assume (G) at worlds $1$ and $j$, (IB) at world~$1$ and (S).
\begin{enumerate}[label=(E\arabic*),leftmargin=*]
\item For every $c>0$ and every non-null $Q\subseteq D\cap\{t>c\}$ there is
$a>0$ with $\Band_j(0,a)\preceq_cQ$.
\item If world~$1$ is complete, then for every $c<\infty$ and every non-null
$A\subseteq D\cap\{t<c\}$ there is $b<\infty$ with
$\Band_j(b,\infty)\prec_dA$.
\end{enumerate}
\end{lemma}
\begin{proof}
(E1) Let $q_0$ and $q_m$ be as in Lemma~\ref{lem:ends}(c). Then
$\alpha_m\coloneqq\alpha_{q_m}\to0$, and $\alpha_m>0$ since $C_{q_m}\ne D_1$.
Let $F$ be $D_1$ silenced on $Q$. Credit at $1$ is debt at $j$ on $D$, so
$F\prec_jD_1=C_{q_0}$, and (T2) gives $C_{q_m}\preceq_1F$ for large $m$. Fix
such $m$ with $\alpha_m<c$, put $R\coloneqq\Band_j(0,\alpha_m)$, and let $N$
be $D_1$ silenced on $Q\cup R$. Then $F=N+\cred(R)$, while $C_{q_m}$ gives
credit on $Q$ and silence or debt on $R$, so $N+\cred(Q)\preceq_1C_{q_m}$ by
(G). Hence $N+\cred(Q)\preceq_1N+\cred(R)$, and by (IB) $R\preceq_cQ$. Take
$a\coloneqq\alpha_m$.

(E2) Let $q_e$ and $q_m\uparrow q_e$ be as in Lemma~\ref{lem:ends}(b), so
$\beta_m\coloneqq\beta_{q_m}\to\infty$. Let $Y$ be $D_j$ with its debt on $A$
silenced; $Y\prec_1D_j=C_{q_e}$ by strict (G), so (T1) gives
$Y\prec_1C_{q_m}$ for large $m$. Fix such $m$ with $\beta_m>c$ and let $N$ be
$D_j$ silenced on $A\cup\Band_j[\beta_m,\infty)$. Then
$Y=N+\debt(\Band_j[\beta_m,\infty))$ and
$C_{q_m}=N+\debt(A)+\cred(\Band_j(\alpha_m,\infty))\preceq_1N+\debt(A)$. So
$N+\debt(\Band_j[\beta_m,\infty))\prec_1N+\debt(A)$, and (IB) gives the claim
with $b\coloneqq\beta_m$.
\end{proof}

\begin{lemma}\label{lem:yard}
Let $n\ge3$, assume (G) at every world, (IB) at world~$1$ and (S), and
fix $j\ge2$.
\begin{enumerate}[label=(\alph*)]
\item If world~$1$ is complete, then for every non-null $A\subseteq H'$ there
are $k\notin\{1,j\}$, a non-null $A_0\subseteq D_{1j}\cap D_{1k}$ with
$A_0\preceq_dA$, and $b<\infty$ with $\Band_j(b,\infty)\prec_dA_0$ and
$\Band_k(b,\infty)\prec_dA_0$. Hence every measurable subset of these two
tails is $\preceq_dA$, and for every open interval $I\subseteq(0,\infty)$
and every $b'\ge b$ the set
\[
 Z_I(A)\coloneqq D_{1j}\cap D_{1k}\cap\{t_k>b'\}\cap\{t_j\in I\}
\]
is non-null and $\preceq_dA$.
\item For every non-null $T\subseteq H'$ there are $k\notin\{1,j\}$ and
$a>0$ with $\Band_j(0,a)\preceq_cT$ and $\Band_k(0,a)\preceq_cT$. Hence for
every open interval $I$ and every $a'\in(0,a]$ the set
\[
 Z^c_I(T)\coloneqq D_{1j}\cap D_{1k}\cap\{t_k<a'\}\cap\{t_j\in I\}
\]
is non-null and $\preceq_cT$.
\end{enumerate}
For $n=3$, $k$ is the third world and $D_{1j}\cap D_{1k}=P_1$. We call
$Z_I(A)$ a \emph{debt yardstick} for $A$ and $Z^c_I(T)$ a \emph{credit
yardstick} for $T$, at levels $I$. The parameters $b'$, $a'$ may be chosen
as needed.
\end{lemma}
\begin{proof}
(a) Since $H'=\bigcup_{l\ge2}D_{1l}$, $A$ meets some $D_{1l}$ in a non-null
set. Put $k\coloneqq l$ if $l\ne j$, and let $k\notin\{1,j\}$ be arbitrary if
$l=j$. Choose $c$ with $A\cap D_{1l}\cap\{t_l<c\}$ non-null; (E2) on
$[e_1,e_l]$ gives $b_0$ with $\Band_l(b_0,\infty)\prec_d
A\cap D_{1l}\cap\{t_l<c\}\preceq_dA$. Let $l'$ be the other element of
$\{j,k\}$ and put
$A_0\coloneqq D_{1j}\cap D_{1k}\cap\{b_0<t_l<b_0+1\}\cap\{t_{l'}<1\}$, which
is non-null by the positive joint density of $(t_j,t_k)$, and
$A_0\preceq_dA$ by (G). Since $A_0$ is bounded in $t_j$ and in $t_k$, (E2)
on $[e_1,e_j]$ and on $[e_1,e_k]$ gives the two tails; take the larger
level. Subsets of a tail are $\preceq_d$ the tail by (G), and $Z_I(A)$ lies
in the $t_k$-tail.

(b) Let $T$ meet $D_{1l}$ in a non-null set, and choose $k$ as in (a). If
$l=k$, choose $c>0$ with $T\cap D_{1k}\cap\{t_k>c\}$ non-null; (E1) gives
$a_k$ with $\Band_k(0,a_k)\preceq_cT$. The set
$T_1\coloneqq D_{1j}\cap D_{1k}\cap\{t_k<a_k\}\cap\{1<t_j<2\}$ is non-null,
$T_1\preceq_cT$ by (G), and (E1) on $[e_1,e_j]$ applied to $T_1$ gives $a_j$
with $\Band_j(0,a_j)\preceq_cT_1$. If $l=j$, exchange the roles of $j$ and
$k$ in this argument. Put $a\coloneqq\min(a_j,a_k)$.
\end{proof}

\begin{corollary}\label{cor:X5}
Let $n\ge3$ and assume (G) at every world, (IB) at world~$1$ and (S).
On every edge $[e_1,e_j]$, the germ point $(0,+)$ is never $T$-heavy, and,
if world~$1$ is complete, $(\infty,-)$ is never $A$-heavy.
\end{corollary}
\begin{proof}
By Lemma~\ref{lem:yard}, $\Band_j(0,a)\preceq_cT$ and
$\Band_j(b,\infty)\preceq_dA$ for some $a>0$ and $b<\infty$.
\end{proof}

\begin{lemma}\label{lem:fineregion}
Let $n\ge3$, assume (G) at every world and (S), and let world $v$ be complete
with (IB). For $w\ne v$, world $v$'s debt order is fine on $D_{vw}$ if and
only if $D_{vw}\in\Jid^v_A$ for every non-null $A\subseteq H'$; likewise for
the credit order and $\Kid^v_T$. World $v$ is fine if and only if it is fine
on $D_{vw}$ for every $w\ne v$.
\end{lemma}
\begin{proof}
One direction is trivial. Conversely, given $A\subseteq H'$,
Lemma~\ref{lem:yard}(a) in world $v$'s frame on $[e_v,e_w]$ gives
$A_0\preceq^v_dA$ with $A_0\subseteq D_{vw}$, so
$D_{vw}\in\Jid^v_{A_0}\subseteq\Jid^v_A$. For credit use the credit
yardstick $Z^c_I(T)\subseteq D_{vw}$ of Lemma~\ref{lem:yard}(b). The last
statement follows since $H'=\bigcup_{w\ne v}D_{vw}$ and the ideals are
closed under finite unions and subsets.
\end{proof}

\subsection{Small AC pairs at the far end of an edge}

\begin{lemma}\label{lem:SAC}
Let world~$1$ be complete with (G) and (IB), and assume (G) at world~$j$ and
(S). Let $C_{q_e}=D_j$ and let $q_m\to q_e$ along $[e_1,e_j]$ with
$\alpha_m\coloneqq\alpha_{q_m}<\infty$ for all $m$. Then for every $c<\infty$
and every non-null $A'\subseteq D\cap\{t<c\}$ there is $m_0$ such that for
all $m\ge m_0$ and every non-null $W\subseteq\Band_j(\alpha_m,\infty)$ the
pair $(A',W)$ is strictly AC.
\end{lemma}
\begin{proof}
Let $Y$ be $D_j$ with its debt on $A'$ silenced. Then $Y\prec_1C_{q_e}$, so
by (T1) $Y\prec_1C_{q_m}$ for large $m$; also $\alpha_m\ge\beta_m\to\infty$
by Lemma~\ref{lem:ends}(a). For $m$ large with $\beta_m>c$ let $N$ be $D_j$
silenced on $A'\cup\Band_j[\beta_m,\infty)$. Then
$N\preceq_1N+\debt(\Band_j[\beta_m,\infty))=Y\prec_1C_{q_m}
=N+\debt(A')+\cred(\Band_j(\alpha_m,\infty))$. By (IB),
$V\prec_1V+\debt(A')+\cred(\Band_j(\alpha_m,\infty))$, and removing credit
from part of the tail makes the right side worse by (G).
\end{proof}

\subsection{SI and AC pairs always exist}

\begin{theorem}\label{thm:SIAC}
Assume (G) at worlds $1$ and $j$ and (S). Neither of the following holds:
\begin{enumerate}[leftmargin=*]
\item[(DF)] for every regular $N$, every non-null $A\subseteq D$ and every
measurable $T\subseteq D\setminus A$ such that $N$ is silent on $A\cup T$:
$N\prec_1N+\debt(A)+\cred(T)$;
\item[(CF)] for every regular $N$, every non-null $T\subseteq D$ and every
measurable $A\subseteq D\setminus T$ such that $N$ is silent on $A\cup T$:
$N+\debt(A)+\cred(T)\prec_1N$.
\end{enumerate}
Consequently, if world~$1$ is complete and satisfies (IB), it has an SI pair and an AC pair, each with non-null debt and credit parts contained
in~$D$.
\end{theorem}
\begin{proof}
All edge models on $[e_1,e_j]$ agree off $D$, and so do the regular models
below, which differ from edge models only on subsets of $D$.

\emph{Two consequences.} (DF$'$) Under (DF), if regular $X,Y$ agree off $D$,
$\debt(X)\subseteq\debt(Y)$ and $\debt(Y)\setminus\debt(X)$ is non-null, then
$X\prec_1Y$. Indeed, let $\mathcal D\coloneqq\{X\ne Y\}\subseteq D$ and let
$N$ be $X$ silenced on $\mathcal D$, which is also $Y$ silenced on
$\mathcal D$. On $\mathcal D$, $X$ never gives debt, so $X\preceq_1N$ by (G).
On $\mathcal D$, $Y$ gives debt exactly on $A\coloneqq\debt(Y)\setminus
\debt(X)$, so $Y=N+\debt(A)+\cred(T)$ with $T\coloneqq\cred(Y)\cap\mathcal D$,
and (DF) gives $N\prec_1Y$. (CF$'$) Under (CF), if $X,Y$ agree off $D$,
$\cred(Y)\subseteq\cred(X)$ and $\cred(X)\setminus\cred(Y)$ is non-null, then
$X\prec_1Y$: let $N$ be $Y$ silenced on $\mathcal D$; on $\mathcal D$, $Y$
never gives credit, so $N\preceq_1Y$; and $X=N+\debt(A)+\cred(T)$ with
$T\coloneqq\cred(X)\cap\mathcal D\supseteq\cred(X)\setminus\cred(Y)$ non-null
and $A\coloneqq\debt(X)\cap\mathcal D$, so $X\prec_1N$ by (CF).

\emph{Suppose (DF).} Step~1: $q\mapsto\beta_q\in[0,\infty]$ is continuous.
If $\beta_q<b<\infty$, then $C_q\prec_1E_j(b)$ by (DF$'$), so
$q\in A_1(E_j(b))$. If $q_m\to p$ and $\lim\beta_{q_m}<\beta_p$, choose $b$
strictly between; closedness gives $C_p\preceq_1E_j(b)$, but
$E_j(b)\prec_1C_p$ by (DF$'$). So $\beta$ is lower semicontinuous. Let
$\beta_p<\infty$ and suppose $q_m\to p$ with $\beta_{q_m}\ge\beta_p+\delta$,
$\delta>0$. Then $p\ne e_j$, and $\beta_{q_m}>0$, so $q_m$ is in the open
edge for large $m$.
Choose a non-null $A\subseteq\Band_j(\beta_p,\beta_p+\delta)$ whose
complement in that band is non-null, and put $F\coloneqq C_p+\debt(A)$;
$C_p$ gives silence or credit on $A$, so $F\prec_jC_p$. By (T2),
$C_{q_m}\preceq_1F$ for large $m$. But $\debt(F)\subseteq
\Band_j[0,\beta_p+\delta)\subseteq\debt(C_{q_m})$ with non-null difference,
so $F\prec_1C_{q_m}$ by (DF$'$). Hence $\beta$ is continuous;
$\beta(e_1)=0$ and $\beta(e_j)=\infty$, so every value in $[0,\infty]$ is
attained.

Step~2: every $p$ with $0<\beta_p<\infty$ is a local extremum of $\beta$.
Such $p$ lies in the open edge. If $\beta_p<\alpha_p$, take a non-null
$Z\subseteq\Band_j(\beta_p,\alpha_p)$ and $E\coloneqq C_p+\cred(Z)$; then
$E\prec_1C_p$, so near $p$, $C_q\not\preceq_1E$. If $\beta_q<\beta_p$, then
$\debt(C_q)\subsetneq\debt(E)=\debt(C_p)$ with non-null difference, and
(DF$'$) gives $C_q\prec_1E$: a contradiction. So $p$ is a local minimum. If
$\beta_p=\alpha_p$, take a non-null $Z\subseteq\Band_j(\alpha_p,\infty)$ and
let $F$ be $C_p$ silenced on $Z$; then $F\prec_jC_p$, and by (T2)
$C_q\preceq_1F$ near $p$. If $\beta_q>\beta_p$, then
$\debt(F)=\debt(C_p)\subsetneq\debt(C_q)$, and (DF$'$) gives $F\prec_1C_q$.
So $p$ is a local maximum.

Step~3. The values of a real function at its local minima on a
second-countable space form a countable set: assign to each such value a
basic open set on which it is the minimum; distinct values get distinct
sets. The same holds for maxima. By Steps~1 and~2 every value in
$(0,\infty)$ is attained, and only at local extrema. So $(0,\infty)$ is
countable, a contradiction.

\emph{Suppose (CF).} The same argument with $\alpha$ in place of $\beta$ and
(CF$'$) in place of (DF$'$). Lower semicontinuity: if $\alpha_q<a<\infty$,
then $C_q\prec_1E^j_1(a)$ by (CF$'$), and $\alpha_p>a$ gives
$E^j_1(a)\prec_1C_p$. Upper semicontinuity at $\alpha_p<\infty$: let $F$ be
$C_p$ with its credit silenced on a non-null proper part $A$ of
$\Band_j(\alpha_p,\alpha_p+\delta)$, where $q_m\to p$ and
$\alpha_{q_m}\ge\alpha_p+\delta$; then $F\prec_jC_p$, the $q_m$ lie in the
open edge for large $m$, (T2) gives $C_{q_m}\preceq_1F$, and
$\cred(C_{q_m})\subsetneq\cred(F)$ with non-null difference gives
$F\prec_1C_{q_m}$. Every $p$ with $0<\alpha_p<\infty$ has $\beta_p>0$ or
$\beta_p<\alpha_p$. If $\beta_p>0$, let $E$ be $C_p$ with its debt silenced
on a non-null $Z\subseteq\Band_j(0,\beta_p)$: $E\prec_1C_p$ and
$\cred(E)=\cred(C_p)$, so $\alpha_q<\alpha_p$ near $p$ would give
$C_q\prec_1E$; $p$ is a local minimum. If $\beta_p<\alpha_p$, let
$F\coloneqq C_p+\debt(Z)$ with a non-null $Z\subseteq\Band_j(\beta_p,\alpha_p)$:
$F\prec_jC_p$ and $\cred(F)=\cred(C_p)$, so $\alpha_q>\alpha_p$ near $p$
would give $F\prec_1C_q$; $p$ is a local maximum. Step~3 concludes.

\emph{Consequence.} Since (DF) fails, there are $N$, a non-null
$A\subseteq D$ and $T\subseteq D\setminus A$ with $N$ silent on $A\cup T$
and $N\not\prec_1N+\debt(A)+\cred(T)$. By completeness and (IB),
$V+\debt(A)+\cred(T)\preceq_1V$. Split $A$ into non-null $A_1,A_2$; by
strict (G), $V+\debt(A_1)+\cred(T)\prec_1V$, and $T$ is non-null since an SI
pair has non-null credit part. Since (CF) fails, there are $N$, a non-null
$T\subseteq D$ and $A\subseteq D\setminus T$ with
$V\preceq_1V+\debt(A)+\cred(T)$; $A$ is non-null by strict (G).
\end{proof}

\subsection{Local extensions from fineness}

\begin{proposition}\label{prop:localext}
Let world $v$ be complete with (G) and (IB), and let $R\subseteq H'$ be
measurable and non-null. Suppose world $v$ is fine on $R$, and some SI pair and some AC pair of world $v$ have non-null debt and credit parts
contained in $R$. Then there are $\kappa\in(0,\infty)$ and finitely additive
probabilities $P_d,P_c$ on the measurable subsets of $R$ with the following
properties.
\begin{enumerate}[label=(\roman*)]
\item For $B,C\subseteq R$: $B\preceq^v_dC$ implies $P_d(B)\le P_d(C)$, and
$P_d(B)<P_d(C)$ implies $B\prec^v_dC$; likewise for $P_c$ and
$\preceq^v_c$.
\item $P_d(B)=0$ if and only if $B$ is null; likewise for $P_c$.
\item Every measurable $B\subseteq R$ has subsets whose $P_d$-mass, resp.\
$P_c$-mass, is arbitrarily close to any prescribed value in $[0,P_d(B)]$,
resp.\ $[0,P_c(B)]$. For every $\eta>0$, $R$ is a finite disjoint union of
measurable sets $Q$ with $P_d(Q)<\eta$ and $P_c(Q)<\eta$.
\item For regular $X,Y$ that agree off $R$, $X\preceq_vY$ implies
$\Lambda(X)\le\Lambda(Y)$, where
\[
 \Lambda(M)\coloneqq P_d(\debt_v(M)\cap R)-\kappa P_c(\cred_v(M)\cap R);
\]
and $\Lambda(X)<\Lambda(Y)$ implies $X\prec_vY$.
\end{enumerate}
If $R=H'$, world $v$ has a loss of CBOI's form satisfying
\eqref{eq:positive} that weakly extends $\preceq_v$ on all regular pairs.
\end{proposition}
\begin{proof}
This is CCF's Sections~5 and~6 and Theorem~15, with $H'$ replaced by $R$. The proof of CCF's
Lemma~8 shows that $\preceq^v_d$ and $\preceq^v_c$, restricted to the
measurable subsets of $R$, are atomless comparative probabilities in which
exactly the non-null sets are $\succ\varnothing$; they are fine by
hypothesis (a finite cover by sets $\preceq A$ becomes a partition by (G)).
CCF's Theorem~9 (Savage--Wakker for algebras,
\cite{Savage,Wakker}) gives the almost agreeing
$P_d,P_c$, and the proofs of CCF's Lemma~10(i)--(iv) are internal to the
algebra. This gives (i)--(iii); the partition in (iii) comes from that proof:
repeated halving gives a non-null $C\subseteq R$ with $P_d(C)<\eta$, a
partition of $R$ into sets $\preceq^v_dC$ has pieces of $P_d$-mass at most
$P_d(C)$ by (i), and a common refinement with the analogous partition for
$P_c$ works for both. CCF's Lemma~11 is the only place where CCF uses (EP)
after fineness; it is replaced by the hypothesis: with the ratios
$r(A,T)=P_d(A)/P_c(T)$ and $\sigma_\pm$ defined as in CCF's Section~6 over
pairs contained in $R$, the given SI pair has ratio in $(0,\infty)$ by (ii),
so $\sigma_->0$, and the given AC pair has finite ratio, so
$\sigma_+<\infty$. CCF's proof of Lemma~12 ($\sigma_-\le\sigma_+$) chooses
its pieces inside the given pairs and $T^*$ inside $R\setminus(A'\cup A_0')$,
whose $P_c$-mass is at least $1-1/K-1/K_0-2\gamma$ because $P_c(R)=1$; all
its comparisons are made in contexts silent on the changed sets. Fix
$\kappa\in[\sigma_-,\sigma_+]$. CCF's Lemma~13 then holds for disjoint
$A,T\subseteq R$, and CCF's proof of Theorem~14, applied to $M\coloneqq X$
and $M'\coloneqq Y$, uses only subsets of $\{X\ne Y\}\subseteq R$; it gives
the first half of (iv). The second half follows from the first and
completeness. For $R=H'$, CCF's Theorem~15 builds the loss from $\Lambda$.
\end{proof}

\begin{corollary}\label{cor:ext}
Let $n\ge3$, assume (G) at every world and (S), and let world $v$ be complete
with (IB).
\begin{enumerate}[label=(\alph*)]
\item If world $v$ is fine on $D_{vw}$, Proposition~\ref{prop:localext}
applies with $R\coloneqq D_{vw}$. We call the resulting $\Lambda$ the
\emph{local extension} of world $v$ on $D_{vw}$, and write it
$\Lambda_v=P^v_d(\debt_v\cap D_{vw})-\kappa_vP^v_c(\cred_v\cap D_{vw})$.
\item If world $v$ is fine, it has a loss of CBOI's form satisfying
\eqref{eq:positive} that weakly extends $\preceq_v$ on all regular pairs.
\end{enumerate}
\end{corollary}
\begin{proof}
Theorem~\ref{thm:SIAC} in world $v$'s frame on $[e_v,e_w]$ supplies the
pairs inside $D_{vw}\subseteq H'$.
\end{proof}

Here \eqref{eq:positive} is CBOI's positivity condition: the charges vanish
exactly on null sets. In the setting of CCF's Theorem~15 (all edges at $v$
with (EP)), Corollary~\ref{cor:ext}(b) is CCF's Theorem~15; the point is that
fineness alone suffices.

\subsection{Transfer at silent selections}

\begin{lemma}\label{lem:silent}
Assume (IB) at worlds $1$ and $j$. Let $C$ be a selection on the open edge
$(e_1,e_j)$, and let $B\subseteq\debt(C)$ and $B'\subseteq\sil(C)$ be
measurable. If $B'\prec_dB$, then $B'\preceq^j_cB$.
\end{lemma}
\begin{proof}
Let $N$ be $C$ silenced on $B\cup B'$, so $C=N+\debt(B)$; put
$X\coloneqq N+\debt(B')$. By (IB), $X\prec_1C$, so $C\not\preceq_1X$, and
competitiveness of $C$ on $\{1,j\}$ gives $C\preceq_jX$. By \eqref{eq:swap},
$C=N+\cred_j(B)$ and $X=N+\cred_j(B')$, and (IB) gives
$V+\cred_j(B)\preceq_jV+\cred_j(B')$, that is, $B'\preceq^j_cB$.
\end{proof}

\section{Edge transfer and one fine world}\label{sec:transfer}

\emph{Standing for Sections~\ref{sec:partner}--\ref{sec:edgefine}:} $n\ge3$;
(G) at every world; (S); worlds $1$ and $j\ne1$ are complete with (IB).
$D\coloneqq D_{1j}$, $t\coloneqq t_j$, and the third index enters only through
the yardsticks of Lemma~\ref{lem:yard}. Nothing is assumed about $V$.

\subsection{Partner bounds}\label{sec:partner}

\begin{definition}\label{def:partner}
A \emph{partner bound} of world $j$ on $D$ is a pair of finite nonnegative
finitely additive charges $a,b$ on the measurable subsets of $D$, vanishing on
null sets and positive on non-null sets, such that for every selection $C$ on
$[e_1,e_j]$ and every regular $X$ that agrees with $C$ off $D$,
\begin{equation}
 C\preceq_jX\ \Longrightarrow\ \Lambda_j(C)\le\Lambda_j(X),\qquad
 \Lambda_j(M)\coloneqq a(D\setminus\debt(M))+b(D\cap\cred(M)).
 \label{eq:partner}
\end{equation}
For an ultrafilter $\Uf$ containing $D$ put
$a(\Uf)\coloneqq\inf\{a(B):B\in\Uf,\ B\subseteq D\}$, and similarly
$b(\Uf)$. $\Uf$ is an \emph{atom} of $a$ if $a(\Uf)>0$.
\end{definition}

By \eqref{eq:swap}, on $D$ world $j$'s rank is $2$ minus world $1$'s, so
$D\setminus\debt(M)=D\cap\{r_j(\cdot,M)\ge1\}$ and
$D\cap\cred(M)=D\cap\{r_j(\cdot,M)\ge2\}$: $\Lambda_j$ is CBOI's form,
restricted to $D$, in world $j$'s ranks. Since world $j$ is complete,
\eqref{eq:partner} gives, for $C$ and $X$ as there,
\begin{equation}
 \Lambda_j(X)<\Lambda_j(C)\ \Longrightarrow\ X\prec_jC.
 \label{eq:partnerstrict}
\end{equation}
If $X$ and $C$ differ exactly on disjoint sets $E_1,\dots,E_r\subseteq D$, on
each of which both have constant advice, $\Lambda_j(X)-\Lambda_j(C)$ is the
sum of the following increments, according to the change of world~$1$'s
advice from $C$ to $X$ on $E_l$:
\begin{equation}
 \debt\to\sil:\ +a(E_l);\qquad\sil\to\cred:\ +b(E_l);\qquad
 \debt\to\cred:\ +a(E_l)+b(E_l);
 \label{eq:incr}
\end{equation}
and the negatives of these for the reverse changes. If $a(\Uf)=0$, then for
every $R\in\Uf$ and $\varepsilon>0$ some $B\in\Uf$ with $B\subseteq R\cap D$
has $a(B)<\varepsilon$; if $a(\Uf)=b(\Uf)=0$, some such $B$ has
$a(B)+b(B)<\varepsilon$ (intersect members of $\Uf$).

\begin{lemma}\label{lem:sources}
\begin{enumerate}[label=(\alph*)]
\item If world $j$ is fine on $D$, its local extension (Corollary~\ref{cor:ext})
gives a partner bound with $a\coloneqq\kappa_jP^j_c$ and $b\coloneqq P^j_d$,
and $a(\Uf)=b(\Uf)=0$ for every ultrafilter $\Uf$ containing $D$.
\item If world $j$ has a loss $L_j$ of CBOI's form with
\eqref{eq:positive} and \eqref{eq:W}, then $a\coloneqq\nu_{j1}|_D$ and
$b\coloneqq\nu_{j2}|_D$ form a partner bound.
\item A strictly positive two-sided local extension of CBOI's form on $D$
(the implication in \eqref{eq:partner} for all regular $X,C$ that agree off
$D$) is a partner bound.
\end{enumerate}
\end{lemma}
\begin{proof}
(a) By \eqref{eq:swap}, $\cred_j=\debt$ and $\debt_j=\cred$ on $D$, so the
local extension is $P^j_d(D\cap\cred(M))-\kappa_jP^j_c(D\cap\debt(M))
=\Lambda_j(M)-\kappa_j$. It weakly extends $\preceq_j$ on all regular pairs
agreeing off $D$, and $a,b$ are positive on non-null sets by
Proposition~\ref{prop:localext}(ii). Given $\eta>0$, partition $D$ into
finitely many sets of $P^j_d$- and $P^j_c$-mass $<\eta$
(Proposition~\ref{prop:localext}(iii)); one of them is in $\Uf$. (b)
Selections are coherent and undominated, and for regular models only
$\nu_{j1},\nu_{j2}$ on $H'$ matter. If $X$ and $C$ agree off $D$, finite
additivity gives $L_j(X)-L_j(C)=\Lambda_j(X)-\Lambda_j(C)$. (c) is clear.
\end{proof}

\subsection{Transfer from a partner bound}

\begin{lemma}\label{lem:creditconf}
Let world $j$ have a partner bound on $D$. Let $\Uf$ be an ultrafilter
containing $D$, avoiding $\Kid_T$ for some non-null $T$, with
$a(\Uf)=b(\Uf)=0$. Then $H'\setminus\cred(C)\notin\Uf$ for every selection
$C$ on the open edge $(e_1,e_j)$ with $\alpha_C<\infty$.
\end{lemma}
\begin{proof}
Edge models give credit off $D$, so $N_C\coloneqq H'\setminus\cred(C)
=D\cap\{t\le\alpha_C\}$. Suppose $N_C\in\Uf$. Let
$B'\coloneqq Z^c_I(T)$ with $I\subseteq(\alpha_C,\infty)$ bounded: then
$B'\subseteq\cred(C)\cap D$, $B'\preceq_cT$ and $b(B')>0$. Choose
$B\in\Uf$, $B\subseteq N_C$, with $a(B)+b(B)<b(B')$; replacing $B$ by
$B\cap\sil(C)$ or $B\cap\debt(C)$, one of which is in $\Uf$, $C$'s advice on
$B$ is constant. By Lemma~\ref{lem:uf}(c), $T\prec_cB$, so $B'\prec_cB$. Let
$N$ be $C$ silenced on $B\cup B'$ and $X\coloneqq N+\cred(B)$. Then
$X\prec_1N+\cred(B')$ by (IB), and $N+\cred(B')\preceq_1C$ by (G), since $C$
is silent or in debt on $B$. So $C\not\preceq_1X$, and competitiveness of $C$
on $\{1,j\}$ gives $C\preceq_jX$. $X$ agrees with $C$ off $B\cup B'$, so
$\Lambda_j(C)\le\Lambda_j(X)$ by \eqref{eq:partner}. But by \eqref{eq:incr}
$\Lambda_j(X)-\Lambda_j(C)\le a(B)+b(B)-b(B')<0$.
\end{proof}

\begin{lemma}\label{lem:SIpartner}
Let world $j$ have a partner bound on $D$, let $T\subseteq H'$ be non-null,
and let $J\subseteq(0,\infty)$ be a bounded open interval: $J\subseteq(1,2)$
if (EP$_{1j}$) holds, and $J\subseteq(\alpha_{C_0},\infty)$ for a fixed
two-sided selection $C_0$ on $[e_1,e_j]$ if (EP$_{1j}$) fails. Then
$T_y\coloneqq Z^c_J(T)$ satisfies $T_y\preceq_cT$ and $T_y\subseteq
D\cap\{t\in J\}$, and some non-null $D'\subseteq D\cap\{t<\inf J\}$ makes
$(D',T_y)$ an SI pair. If (EP$_{1j}$) holds, the partner bound is not
used.
\end{lemma}
\begin{proof}
$T_y$ is given by Lemma~\ref{lem:yard}(b). Suppose (EP$_{1j}$). By CCF's
Lemmas~3 and~4 (with $s^*=1$ approached from below) there is $q^*$ with $C_{q^*}=E_j(1)$, every neighbourhood of which
contains indices $q$ with $C_q=E_j(s)$ for all $s$ in some interval
$[s_0,1)$. Let $Y\coloneqq E_j(1)+\cred(T_y)$; $E_j(1)$ is silent on
$T_y\subseteq\Band_j(1,2)$, so $Y\prec_1E_j(1)$, and (T1) gives
$Y\prec_1E_j(s)$ for some $s<1$. Since
$Y=E_j(s)+\debt(\Band_j[s,1))+\cred(T_y)$, the pair $(\Band_j[s,1),T_y)$ is
SI by (IB). This is CCF's Lemma~11 at level~$1$.

Suppose (EP$_{1j}$) fails, and let $C_0=(\alpha_0,\beta_0)$, so
$0<\beta_0\le\alpha_0<\infty$ and $T_y\subseteq\cred(C_0)\cap D$. Suppose every
pair $(D',T_y)$ with $D'\subseteq\debt(C_0)\cap D$ non-null is AC. By FAED's
Lemma~1 there is a non-null $D''\subseteq\debt(C_0)\cap D$ with
$a(D'')<b(T_y)$; split it into non-null $D'$ and $D''\setminus D'$. For
$E\in\{D',D''\}$ let $K_E$ be $C_0$ silenced on $E\cup T_y$. Then
$C_0=K_E+\debt(E)+\cred(T_y)$, so $K_E\preceq_1C_0$ by the AC pair and (IB).
By \eqref{eq:incr}, $\Lambda_j(K_E)-\Lambda_j(C_0)=a(E)-b(T_y)<0$, so
$K_E\prec_jC_0$ by \eqref{eq:partnerstrict}, and competitiveness gives
$C_0\preceq_1K_E$. Hence $K_{D'}\sim_1C_0\sim_1K_{D''}$. But $K_{D''}$ is
$K_{D'}$ with its debt on the non-null $D''\setminus D'$ silenced, so
$K_{D''}\prec_1K_{D'}$ by strict (G): a contradiction. So some non-null
$D'\subseteq\debt(C_0)\cap D=\Band_j(0,\beta_0)$ is an SI partner of
$T_y$, and $\beta_0\le\alpha_0<\inf J$.
\end{proof}

\begin{theorem}\label{thm:transfer}
Let world $j$ have a partner bound $(a,b)$ on $D$, and let $\Uf$ be an
ultrafilter containing $D$.
\begin{enumerate}[label=(\alph*)]
\item If $\Uf$ avoids $\Jid_A$ for some non-null $A\subseteq H'$, then
$a(\Uf)>0$.
\item If $\Uf$ avoids $\Kid_T$ for some non-null $T\subseteq H'$, then
$a(\Uf)>0$ or $b(\Uf)>0$.
\end{enumerate}
So every heavy ultrafilter of world~$1$ on $D$ is an atom of the partner's
charges. Nothing is assumed about $V$, and neither completeness nor (IB) is
used at worlds other than $1$ and $j$.
\end{theorem}
\begin{proof}
\emph{Preliminaries.} By Lemma~\ref{lem:uf}(d), $\Uf$ contains all one-sided
bands of a unique germ point $(z,\sigma)$ of $[e_1,e_j]$, which is $A$-heavy
in~(a) and $T$-heavy in~(b). Let $q_e$ and $q_m\uparrow q_e$ be as in
Lemma~\ref{lem:ends}(b), and $C_m\coloneqq C_{q_m}=(\alpha_m,\beta_m)$:
selections on the open edge, $C_m\ne D_j$, so $\beta_m<\infty$, and
$\beta_m\to\infty$.

(a) Suppose $a(\Uf)=0$. By Corollary~\ref{cor:X5}, $z<\infty$. For
$\beta_m>z+1$, $\debt(C_m)=\Band_j(0,\beta_m)$ contains a one-sided band of
$(z,\sigma)$, so $\debt(C_m)\in\Uf$. Pass to a subsequence along which either
$\beta_m<\alpha_m$ for all $m$ (\emph{silent}) or $\alpha_m=\beta_m$ for all
$m$ (\emph{precise}).

\emph{Silent.} Fix $m$ in the subsequence with $\beta_m>z+1$. Let $B'\coloneqq Z_I(A)$ with $I$
a bounded open interval in $(\beta_m,\alpha_m)$: $B'\subseteq\sil(C_m)$,
$B'\preceq_dA$ and $a(B')>0$. Choose $B\in\Uf$, $B\subseteq\debt(C_m)$, with
$a(B)<a(B')$. By Lemma~\ref{lem:uf}(c), $B'\prec_dB$. Let $N$ be $C_m$
silenced on $B\cup B'$, so $C_m=N+\debt(B)$, and put $X\coloneqq
N+\debt(B')$. By (IB), $X\prec_1C_m$; by competitiveness $C_m\preceq_jX$; by
\eqref{eq:partner} $\Lambda_j(C_m)\le\Lambda_j(X)$. But by \eqref{eq:incr}
$\Lambda_j(X)-\Lambda_j(C_m)=a(B)-a(B')<0$.

\emph{Precise.} Then $\alpha_m=\beta_m\to\infty$. Let $A'\coloneqq Z_I(A)$
with $I\coloneqq(1,2)$, so $A'\subseteq D\cap\{t<2\}$ and $A'\preceq_dA$, and
let $m_0$ be given by Lemma~\ref{lem:SAC} for $A'$ along the subsequence. Fix
$m\ge m_0$ in the subsequence with $\beta_m>z+3$ and then $m'\ge m$ in it
with $\alpha_{m'}>\alpha_m$, and put $W\coloneqq\Band_j(\alpha_{m'},\infty)$.
Then $A'\subseteq\debt(C_m)$, $W\subseteq\cred(C_m)\cap D$, $A'\cap W=\varnothing$, $b(W)>0$, and $(A',W)$
is strictly AC. Since $A'\in\Jid_A$, $A'\notin\Uf$, so
$\debt(C_m)\setminus A'\in\Uf$; choose $B\in\Uf$ inside it with
$a(B)<b(W)$. By Lemma~\ref{lem:uf}(c), $A'\prec_dB$. Let $K$ be $C_m$
silenced on the pairwise disjoint sets $A'$, $B$, $W$. By (IB) and the AC
pair, $K\preceq_1K+\debt(A')+\cred(W)$, and by (IB) and $A'\prec_dB$,
$K+\cred(W)+\debt(A')\prec_1K+\cred(W)+\debt(B)$; so
$K\prec_1K+\cred(W)+\debt(B)$. Let $X$ be $C_m$ with $B$ and $W$ silenced.
The pairs $(X,C_m)$ and $(K,K+\cred(W)+\debt(B))$ each agree off $B\cup W$
and carry the same advice on $B\cup W$, so $X\prec_1C_m$ by (IB). By
competitiveness and \eqref{eq:partner}, $\Lambda_j(C_m)\le\Lambda_j(X)$; but
$\Lambda_j(X)-\Lambda_j(C_m)=a(B)-b(W)<0$.

(b) Suppose $a(\Uf)=b(\Uf)=0$. By Corollary~\ref{cor:X5}, $z>0$; write
$w\coloneqq z\in(0,\infty]$. Fix $J$, $T_y$ and $D'$ as in
Lemma~\ref{lem:SIpartner}; $\Uf$ avoids $\Kid_{T_y}\subseteq\Kid_T$
(Lemma~\ref{lem:ideal}). Fix a non-null $D'_0\subseteq D'$ with
$D'\setminus D'_0$ non-null; then $D'_0\prec_dD'$ by strict (G), so every set
$\preceq_dD'_0$ is $\prec_dD'$.

\emph{Step 1: $0<w<\infty$.} (i) Suppose $\alpha_m<\infty$ for infinitely
many $m$. For such $m$ with $\beta_m>w+1$, the bands of $(w,\sigma)$ of size
$<1$ lie in $D\cap\{t\le\alpha_m\}=H'\setminus\cred(C_m)$, which is
therefore in $\Uf$. This contradicts Lemma~\ref{lem:creditconf}, applied
with $T_y$ in place of $T$.

(ii) Otherwise $C_m=E_j(\beta_m)$ for all large $m$. Fix such $m$ with
$\beta\coloneqq\beta_m>\max(w,\sup J)+1$, and $E\coloneqq C_m$. Let
$Z\coloneqq Z_I(D'_0)$ with $I\coloneqq(\beta,\beta+1)$: $Z\subseteq\sil(E)$,
$Z\prec_dD'$, $a(Z)>0$, and $Z$ misses $T_y$ and $D'$. The set
$\debt(E)=\Band_j(0,\beta)$ contains bands of $(w,\sigma)$, so it is in
$\Uf$; $T_y\in\Kid_{T_y}$ is not; choose $B\in\Uf$,
$B\subseteq\Band_j(0,\beta)\setminus T_y$, with $a(B)+b(B)<a(Z)$. Let $X$
be $E$ with credit instead of debt on $B$ and debt instead of silence on
$Z$. By \eqref{eq:incr}, $\Lambda_j(X)-\Lambda_j(E)=a(B)+b(B)-a(Z)<0$, so
$X\prec_jE$ by \eqref{eq:partnerstrict}, and competitiveness gives
$E\preceq_1X$. With $N$ the model $E$ silenced on $B\cup Z$,
$N\preceq_1N+\debt(B)=E\preceq_1X=N+\cred(B)+\debt(Z)$, so (IB) gives
$V\preceq_1V+\debt(Z)+\cred(B)$. By Lemma~\ref{lem:uf}(c), $T_y\prec_cB$;
in the context $V+\debt(Z)$, which is silent on $B\cup T_y$, this gives
$V+\debt(Z)+\cred(B)\prec_1V+\debt(Z)+\cred(T_y)$. Hence
$V\prec_1V+\debt(Z)+\cred(T_y)$. But $Z\prec_dD'$ in the context
$V+\cred(T_y)$, silent on $Z\cup D'$, gives
$V+\debt(Z)+\cred(T_y)\prec_1V+\debt(D')+\cred(T_y)\prec_1V$.

\emph{Step 2: $w=\infty$.} $\Uf$ contains every tail $\Band_j(M,\infty)$.
(i) Suppose $\alpha_m<\infty$ for infinitely many $m$. Let
$A'\coloneqq Z_I(D'_0)$ with $I$ a bounded open interval in $(0,\inf J)$
and $c\coloneqq\sup I$: $A'\subseteq D\cap\{t<c\}$, $A'\prec_dD'$,
$A'\cap T_y=\varnothing$. By Lemma~\ref{lem:SAC}, for large such $m$ with
$\alpha_m>\sup J$ the pair $(A',W_m)$, $W_m\coloneqq\Band_j(\alpha_m,\infty)$,
is strictly AC. Then
\[
 V+\debt(A')+\cred(T_y)\prec_1V+\debt(D')+\cred(T_y)\prec_1V
 \prec_1V+\debt(A')+\cred(W_m),
\]
the first step in the context $V+\cred(T_y)$, silent on $A'\cup D'$. (IB) in
the context $V+\debt(A')$, silent on $T_y\cup W_m$, gives $W_m\prec_cT_y$.
So this tail is in $\Kid_{T_y}$, and it is in $\Uf$: a contradiction.

(ii) Otherwise $C_m=E_j(\beta_m)$ for all large $m$; fix
$E\coloneqq C_m$ with $\beta\coloneqq\beta_m>\sup J+1$, and $Z$ as in
Step~1(ii). Choose $W\in\Uf$, $W\subseteq\Band_j(\beta+1,\infty)$, with
$b(W)<a(Z)$, and let $X\coloneqq E+\cred(W)+\debt(Z)$; $E$ is silent on the
disjoint sets $W$, $Z$. Then $\Lambda_j(X)-\Lambda_j(E)=b(W)-a(Z)<0$, so
$X\prec_jE$, $E\preceq_1X$, and by (IB) $V\preceq_1V+\debt(Z)+\cred(W)$. As
$T_y\prec_cW$, $V\prec_1V+\debt(Z)+\cred(T_y)$, and $Z\prec_dD'$ gives the
opposite strict inequality as in Step~1(ii).

Every germ point is $(0,+)$, finite, or $(\infty,-)$, so all cases are
covered.
\end{proof}

The proof works at the far end of the edge, where every selection has its
debt band beyond the germ level. Near $q_e$ each selection is either silent
on a band, which hosts a debt yardstick, or has a credit tail, which with the
small AC pair of Lemma~\ref{lem:SAC} pays for removing the heavy debt; no
analysis of how the selections cross the germ level is needed. The partner's
charges enter only through the choice of pieces of small charge in $\Uf$,
which is possible exactly when $\Uf$ is not an atom.

\subsection{Edge transfer of fineness}\label{sec:edgefine}

\begin{corollary}\label{cor:edgefine}
If world $j$ is fine on $D_{1j}$, then $D_{1j}\in\Jid_A\cap\Kid_T$ for all
non-null $A,T\subseteq H'$; in particular world~$1$ is fine on $D_{1j}$. So
fineness on $D_{vw}$ is a property of the edge $[e_v,e_w]$: it holds for
one of its worlds if and only if it holds for the other, provided both are
complete with (IB).
\end{corollary}
\begin{proof}
Lemma~\ref{lem:sources}(a) gives a partner bound with $a(\Uf)=b(\Uf)=0$ for
every ultrafilter $\Uf\ni D$. If $D\notin\Jid_A$, Lemma~\ref{lem:uf}(b) gives
an ultrafilter containing $D$ and avoiding $\Jid_A$, which contradicts
Theorem~\ref{thm:transfer}(a). The credit case uses (b).
\end{proof}

\noindent The partner must be complete with (IB): its debt and credit
orders are comparative probabilities only under (IB), and the strict form
\eqref{eq:partnerstrict} needs completeness. (G) at every world and (S) are
used through the yardsticks of Lemma~\ref{lem:yard}, which is where the third
world enters; for $n=2$ the argument does not run.

\subsection{The third region}

\begin{proposition}\label{prop:third}
Let $n=3$. Let worlds $1$ and $3$ be complete with (G) and (IB), assume (G)
at world~$2$ and (S), and let world $3$ be fine on $D_{32}$, with local
extension $\Lambda_3=P^3_d(\debt_3\cap D_{32})-\kappa_3P^3_c(\cred_3\cap
D_{32})$.
\begin{enumerate}[label=(\alph*)]
\item No ultrafilter that avoids $\Jid_A$, for some non-null $A$, contains
$P_3$.
\item Suppose moreover that world~$1$'s debt order is fine and that
\begin{enumerate}[label=(H$_{\mathrm{SI}}$),leftmargin=*]
\item for every non-null $T$ there are a non-null $T_y\preceq_cT$ with
$T_y\cap P_3=\varnothing$ and a non-null $D'$ such that $(D',T_y)$ is
an SI pair of world~$1$.
\end{enumerate}
Then no ultrafilter that avoids $\Kid_T$, for some non-null $T$, contains
$P_3$.
\end{enumerate}
\end{proposition}
\begin{proof}
Work on $[e_1,e_3]$ with $t\coloneqq t_3$ on $D_{13}\supseteq P_3$. Since
$P_3=D_{13}\cap D_{32}$, models that agree off a subset of $P_3$ agree off
$D_{32}$, so $\Lambda_3$ and its strict form apply to them; and on $P_3$
debt at~$1$ is credit at~$3$ and credit at $1$ is debt at $3$, so the
increments \eqref{eq:incr} hold for changes inside $P_3$ with
$a\coloneqq\kappa_3P^3_c$ and $b\coloneqq P^3_d$. Pieces of small $P^3_d$-
and $P^3_c$-mass exist by Proposition~\ref{prop:localext}(iii). Let $q_e$,
$q_m\uparrow q_e$ and $C_m=(\alpha_m,\beta_m)$ be as in
Lemma~\ref{lem:ends}(b) on $[e_1,e_3]$; $\beta_m\to\infty$.

(a) Let $\Uf\ni P_3$ avoid $\Jid_A$. By Lemma~\ref{lem:uf}(d) it contains
all bands of an $A$-heavy germ point $(z,\sigma)$ of $[e_1,e_3]$, and
$z<\infty$ by Corollary~\ref{cor:X5}. For $\beta_m>z+1$,
$\debt(C_m)\cap P_3\in\Uf$.

\emph{Case $\alpha_m=\infty$ along a subsequence.} Fix such $m$ with
$\beta_m>z+1$. By Lemma~\ref{lem:yard}(a) there is $b_3$ such that every
subset of $\Band_3(b_3,\infty)$ is $\preceq_dA$. The set
$B'\coloneqq P_3\cap\Band_3(\max(b_3,\beta_m),\infty)$ is non-null, lies in
$\sil(C_m)$, and $B'\preceq_dA$. Choose a non-null $B'_0\subseteq B'$ with
$B'\setminus B'_0$ non-null; then $B'_0\prec^3_cB'$ by strict (G) at
world~$3$. Since world~$3$ is fine on $D_{32}\supseteq B'_0$, the set
$\debt(C_m)\cap P_3$ has a subset $B\in\Uf$ with $B\preceq^3_cB'_0$. By
Lemma~\ref{lem:uf}(c), $B'\prec_dB$, and Lemma~\ref{lem:silent} (edge
$[e_1,e_3]$, partner~$3$) gives $B'\preceq^3_cB\preceq^3_cB'_0\prec^3_cB'$:
a contradiction.

\emph{Case $\alpha_m<\infty$ along a subsequence.} Let $A'\coloneqq Z_I(A)$
on $[e_1,e_3]$ with $I\coloneqq(1,2)$; $A'\subseteq P_1\cap\{t<2\}$,
$A'\preceq_dA$. Let $m_0$ be given by Lemma~\ref{lem:SAC}. Fix $m\ge m_0$ in
the subsequence with $\beta_m>z+3$, then $m'\ge m$ in it with
$\alpha_{m'}>\alpha_m$, and put $W\coloneqq P_3\cap\Band_3(\alpha_{m'},
\infty)$: non-null, in $\cred(C_m)$, disjoint from $A'$, and $(A',W)$ is
strictly AC. Choose $B\in\Uf$, $B\subseteq\debt(C_m)\cap P_3$, with
$\kappa_3P^3_c(B)<P^3_d(W)$. Then $A'\prec_dB$, and the argument of the
precise case of Theorem~\ref{thm:transfer}(a) gives $X\prec_1C_m$ for $X$ the
model $C_m$ with $B$ and $W$ silenced. Competitiveness on $\{1,3\}$ gives
$C_m\preceq_3X$, hence $\Lambda_3(C_m)\le\Lambda_3(X)$; but
$\Lambda_3(X)-\Lambda_3(C_m)=\kappa_3P^3_c(B)-P^3_d(W)<0$.

Since $C_m\ne D_3$, $\beta_m<\infty$, so the two cases are exhaustive.

(b) Let $\Uf\ni P_3$ avoid $\Kid_T$. By Lemma~\ref{lem:uf}(d) it contains all
bands of a $T$-heavy germ point $(w,\sigma)$ of $[e_1,e_3]$, and $w>0$. Take
$T_y$, $D'$ from (H$_{\mathrm{SI}}$); $\Uf$ avoids $\Kid_{T_y}$. Fix a
non-null $D'_0\subseteq D'$ with $D'\setminus D'_0$ non-null, so every set
$\preceq_dD'_0$ is $\prec_dD'$. By Lemma~\ref{lem:yard}(a) there is $b$ such
that every subset of $\Band_3(b,\infty)$ is $\preceq_dD'_0$. By fineness of
world~$1$'s debt order, $H'$ is a finite union of sets $\preceq_dD'_0$; one
of them meets $P_3\cap\Band_3(0,1)$ in a non-null set $A'$. Then
$A'\prec_dD'$, $A'\cap T_y=\varnothing$ and $A'\subseteq D_{13}\cap\{t<1\}$.

\emph{Finite $w$, case $\alpha_m<\infty$ along a subsequence.} Let $m_0$ be
given by Lemma~\ref{lem:SAC} for $A'$. Fix $m\ge m_0$ in the subsequence
with $\beta_m>w+1$ and $m'\ge m_0$ in it with $\alpha_{m'}>\alpha_m$, and put
$W\coloneqq P_3\cap\Band_3(\alpha_{m'},\infty)\subseteq\cred(C_m)$. As in
Step~2(i) of Theorem~\ref{thm:transfer}(b),
\[
 V+\debt(A')+\cred(T_y)\prec_1V+\debt(D')+\cred(T_y)\prec_1V
 \prec_1V+\debt(A')+\cred(W),
\]
and (IB) gives $W\prec_cT_y$. Choose $B\in\Uf$,
$B\subseteq\debt(C_m)\cap P_3$, with $P^3_d(B)+\kappa_3P^3_c(B)<P^3_d(W)$.
Since $T_y\prec_cB$, $W\prec_cB$. Let $N$ be $C_m$ silenced on $B\cup W$ and
$X\coloneqq N+\cred(B)$. Then $X\prec_1N+\cred(W)\preceq_1
N+\cred(W)+\debt(B)=C_m$, and
$\Lambda_3(X)-\Lambda_3(C_m)=P^3_d(B)+\kappa_3P^3_c(B)-P^3_d(W)<0$, so
$X\prec_3C_m$. Thus $C_m$ is beaten at both worlds of its support, contrary
to competitiveness.

\emph{Finite $w$, case $\alpha_m=\infty$ along a subsequence.} Fix such $m$
with $\beta\coloneqq\beta_m>w+1$, $E\coloneqq C_m=E_3(\beta)$, and put
$\beta'\coloneqq\max(\beta,b)$, $Z\coloneqq P_3\cap\Band_3(\beta',\beta'+1)$.
Then $Z\subseteq\sil(E)$, $Z\prec_dD'$ and $Z\cap T_y=\varnothing$. Choose
$B\in\Uf$, $B\subseteq P_3\cap\Band_3(0,\beta)$, with
$P^3_d(B)+\kappa_3P^3_c(B)<\kappa_3P^3_c(Z)$, and let $X$ be $E$ with credit
instead of debt on $B$ and debt instead of silence on $Z$. Then
$\Lambda_3(X)-\Lambda_3(E)<0$, so $X\prec_3E$ and $E\preceq_1X$, and the
argument of Step~1(ii) of Theorem~\ref{thm:transfer}(b) gives a
contradiction. (This case does not use the fineness of world~$1$'s debt
order.)

\emph{$w=\infty$.} $\Uf$ contains $P_3\cap\Band_3(M,\infty)$ for every
$M$. If $\alpha_m<\infty$ for infinitely many $m$, the chain of the first
case gives $P_3\cap\Band_3(\alpha_m,\infty)\prec_cT_y$ for large such $m$, so
this member of $\Uf$ is in $\Kid_{T_y}$. Otherwise fix $E\coloneqq
C_m=E_3(\beta)$ and $Z$ as in the second case, choose $W\in\Uf$,
$W\subseteq P_3\cap\Band_3(\beta'+1,\infty)$, with
$P^3_d(W)<\kappa_3P^3_c(Z)$, and let $X\coloneqq E+\cred(W)+\debt(Z)$. Then
$X\prec_3E$, $E\preceq_1X$, $V\preceq_1V+\debt(Z)+\cred(W)$, and
$T_y\prec_cW$ and $Z\prec_dD'$ give a contradiction as in Step~2(ii) of
Theorem~\ref{thm:transfer}(b).
\end{proof}

We do not know whether (H$_{\mathrm{SI}}$) follows from the other
hypotheses of (b): the credit yardsticks of Lemma~\ref{lem:yard}(b) lie in $P_1$, and an SI
partner for them comes from a partner bound on $[e_1,e_2]$ or from (EP) on an
edge at world~$1$ (Lemma~\ref{lem:SIpartner}). It holds in
Theorem~\ref{thm:onefine}.

\subsection{One fine world}

\begin{theorem}\label{thm:onefine}
Let $n=3$, let all three rankings be complete with (G) and (IB), and assume
(S). If one world is fine, then all three are fine. Consequently each world
has a loss of CBOI's form satisfying \eqref{eq:positive} that weakly extends
its ranking on all regular pairs, and every coherent model whose credal set
contains three vertices, $V$ in particular, is dominated by a coherent
model. The first conclusion uses no assumption on~$V$.
\end{theorem}
\begin{proof}
Relabel so that world~$2$ is fine.
\begin{enumerate}[label=\arabic*.,leftmargin=*,itemsep=1pt]
\item Corollary~\ref{cor:edgefine} (base $1$, partner $2$):
$D_{12}\in\Jid_A\cap\Kid_T$ for all non-null $A,T$.
\item Corollary~\ref{cor:edgefine} (base $3$, partner $2$): world~$3$ is
fine on $D_{32}$.
\item World~$1$'s debt order. If $H'\notin\Jid_A$, Lemma~\ref{lem:uf}(b)
gives an ultrafilter avoiding $\Jid_A$. It contains exactly one of
$P_1,P_2,P_3$; not $P_1$ or $P_2$, which lie in $D_{12}\in\Jid_A$; and not
$P_3$, by Proposition~\ref{prop:third}(a). So $H'\in\Jid_A$ for every $A$.
\item World~$1$'s credit order. The same argument with $\Kid_T$, using
Proposition~\ref{prop:third}(b): world~$1$'s debt order is fine by step~3,
and (H$_{\mathrm{SI}}$) holds by Lemma~\ref{lem:SIpartner} on $[e_1,e_2]$
with world~$2$'s partner bound of Lemma~\ref{lem:sources}(a), since
$T_y=Z^c_J(T)\subseteq P_1$. So world~$1$ is fine.
\item World~$3$. World~$1$ is fine on $D_{31}$, so
Corollary~\ref{cor:edgefine} (base $3$, partner $1$) makes world~$3$ fine on
$D_{31}$; with step~2 and Lemma~\ref{lem:fineregion}, world~$3$ is fine.
\item Corollary~\ref{cor:ext}(b) at each world gives the losses; they satisfy
\eqref{eq:W}, and CCF's Lemma~1 gives the dominance.
\end{enumerate}
\end{proof}

\noindent Step~5 avoids a second use of Proposition~\ref{prop:third}. For
the conclusion that $V$ is dominated, Theorem~\ref{thm:fineedge} and
Remark~\ref{rem:edgeroute} below give a second route that avoids
Proposition~\ref{prop:third} altogether: by step~2, the edge $[e_2,e_3]$ is
fine. For $n\ge4$, Corollary~\ref{cor:edgefine} still
transfers fineness along every edge, but we have not checked the analogue of
Proposition~\ref{prop:third}; whether one fine world makes all worlds fine
for $n\ge4$ is open.

\section{One fine edge}\label{sec:edge}

Call the edge $[e_v,e_w]$ \emph{fine} if worlds $v$ and $w$ are fine on
$D_{vw}$. If both are complete with (IB), one of them suffices, by
Corollary~\ref{cor:edgefine}.

\begin{lemma}\label{lem:EPfine}
Let $n\ge3$, assume (G) at every world and (S), and let worlds $v$ and $w$ be
complete with (IB). If (EP$_{vw}$) holds, the edge $[e_v,e_w]$ is fine.
\end{lemma}
\begin{proof}
Take $v=1$, $w=j$. Step~1 of CCF's proof of Theorem~7 and its credit version
use only the edge $[e_1,e_j]$, (EP$_{1j}$), (G) at worlds $1$ and $j$, (S),
and completeness and (IB) of world~$1$. They show: for every non-null
$A_0\subseteq D_{1j}\cap\{t_j\ge a_1\}$ with $a_1>0$, $D_{1j}$ is a finite
union of sets $\preceq_dA_0$, and for every non-null
$T_0\subseteq D_{1j}\cap\{t_j>c\}$ with $c>0$, $D_{1j}$ is a finite union of
sets $\preceq_cT_0$. A non-null $A\subseteq D_{1j}$ contains such an $A_0$,
and $A_0\preceq_dA$ by (G); likewise for credit. So world~$1$ is fine on
$D_{1j}$. (EP$_{1j}$) is symmetric in the two worlds of the edge, so the same
argument in world $j$'s frame applies to world~$j$.
\end{proof}

For a selection $E=(\alpha_E,\beta_E)$ on $[e_1,e_2]$, in world~$1$'s frame,
put
\[
 T_0(E)\coloneqq P_2\cap\{t_2>\alpha_E\},\qquad
 B_0(E)\coloneqq P_1\cap\{t_2<\beta_E\}.
\]
$E$ gives $o_{13}$ (credit at $1$ and $3$, debt at $2$) on $T_0(E)$ and
$o_{23}$ (debt at $1$, credit at $2$ and $3$) on $B_0(E)$. For
$T\subseteq T_0(E)$ and $B\subseteq B_0(E)$ put
$\tau_{T,B}\coloneqq o_{13}(T)+o_{23}(B)$, and write $V+\tau_{T,B}$ for $V$
with these options on $T$ and $B$.

\begin{lemma}\label{lem:XOR}
Assume (IB) at worlds $1$ and $2$. Let $E$ be a selection on the open edge
$(e_1,e_2)$, $T\subseteq T_0(E)$ and $B\subseteq B_0(E)$. Then
$V+\tau_{T,B}\preceq_1V$ or $V+\tau_{T,B}\preceq_2V$. If moreover (G) holds
at world~$3$, $V$ is undominated and $T\cup B$ is non-null, then not both.
\end{lemma}
\begin{proof}
The pairs $(E,E')$, with $E'$ the model $E$ silenced on $T\cup B$, and
$(V+\tau_{T,B},V)$ differ exactly on $T\cup B$, with the same advice there.
By (IB), $E\preceq_kE'$ if and only if $V+\tau_{T,B}\preceq_kV$, for
$k=1,2$; and $E$ is competitive on $\{1,2\}$. For the second statement:
$o_{13}$ and $o_{23}$ are credit at world~$3$, so $V+\tau_{T,B}\prec_3V$ by
strict (G), and if both weak inequalities held, $V+\tau_{T,B}$ would dominate
$V$.
\end{proof}

\emph{Standing for the rest of this section:} $n=3$; all three rankings are
complete with (G) and (IB); (S); and the edge $[e_2,e_3]$ is fine, with local
extensions $\Lambda_2,\Lambda_3$ on $D_{23}$ (Corollary~\ref{cor:ext}) and
probabilities $P^2_d,P^2_c,P^3_d,P^3_c$ and rates $\kappa_2,\kappa_3$. Nothing
is assumed about world~$1$ beyond completeness, (G) and (IB).

\begin{lemma}\label{lem:B1}
For every non-null $B\subseteq P_1$ there is a non-null $Z_B\subseteq P_3$
such that $V+o_{12}(Z)+o_{23}(B)\prec_3V$ for every non-null
$Z\subseteq Z_B$. Only world~$3$'s fineness on $D_{23}$ is used, not
world~$2$'s.
\end{lemma}
\begin{proof}
We work in world~$3$'s frame, with $s\coloneqq t^3_2$ on $D_{32}=D_{23}$. On
$P_3$, $o_{12}$ is debt at $3$; on $P_1$, $o_{23}$ is credit at $3$; so the
claim is that $(Z,B)$ is an SI pair of world~$3$.

(1) By Lemma~\ref{lem:yard}(b) at world~$3$ there is $a'>0$ with
$L\coloneqq\Lo32{a'}\preceq^3_cB$. (2) Choose a non-null $T''\subseteq L$
with $D_{23}\setminus T''$ non-null, and then (Proposition~\ref{prop:localext}
(ii), (iii)) a non-null $A\subseteq D_{23}\setminus T''$ with
$P^3_d(A)<\kappa_3P^3_c(T'')$. The models $V+\debt_3(A)+\cred_3(T'')$ and $V$
agree off $D_{23}$, and $\Lambda_3$ decreases by
$\kappa_3P^3_c(T'')-P^3_d(A)>0$, so $V+\debt_3(A)+\cred_3(T'')\prec_3V$.
(3) By Lemma~\ref{lem:yard}(a) at world~$3$ there is $b$ with
$\Hi32b\prec^3_dA_0\preceq^3_dA$. Put
$Z_B\coloneqq P_3\cap\Hi32{\max(b,a')}$; every non-null $Z\subseteq Z_B$ has
$Z\prec^3_dA$. (4) $Z$ and $T''$ are disjoint (levels of $s$), $A\cap
T''=\varnothing$, and $B\subseteq P_1$ misses $Z,T'',A$. Hence
\[
 V+\debt_3(Z)+\cred_3(B)\preceq_3V+\debt_3(Z)+\cred_3(T'')
 \prec_3V+\debt_3(A)+\cred_3(T'')\prec_3V.
\]
The first step is $T''\preceq^3_cL\preceq^3_cB$ in the context
$V+\debt_3(Z)$, silent on $B\cup T''$; the second is $Z\prec^3_dA$ in the
context $V+\cred_3(T'')$, silent on $Z\cup A$; the third is step~(2).
\end{proof}

\begin{lemma}\label{lem:B2}
Assume moreover that $V$ is undominated. Let $E$ be a selection on the open
edge $(e_1,e_2)$ and $B\subseteq B_0(E)$ non-null. Then
$V\prec_1V+\tau_{T,B}$ for every $T\subseteq T_0(E)$.
\end{lemma}
\begin{proof}
Suppose $V+\tau_{T^*,B}\preceq_1V$ for some $T^*\subseteq T_0(E)$. Let
$Z\coloneqq Z_B$ (Lemma~\ref{lem:B1}) and $\eta\coloneqq\kappa_2P^2_c(Z)>0$.
Write $T^*=Q_1\cup\dots\cup Q_r$, disjoint, with $P^2_d(Q_k)<\eta$
(Proposition~\ref{prop:localext}(iii); $T^*\subseteq P_2\subseteq D_{23}$),
and $T_k\coloneqq Q_1\cup\dots\cup Q_k$. Since $V+\tau_{\varnothing,B}=
V+o_{23}(B)$ has debt at $1$ on $B$, $V\prec_1V+\tau_{\varnothing,B}$; so
there is a least $k^*\ge1$ with $V+\tau_{T_{k^*},B}\preceq_1V$, and
$V+\tau_{T_{k^*-1},B}\preceq_2V$ by Lemma~\ref{lem:XOR}. Put
\[
 d\coloneqq o_{13}(T_{k^*})+o_{23}(B)+o_{12}(Z),
\]
a change on the disjoint sets $T_{k^*}\subseteq P_2$, $B\subseteq P_1$,
$Z\subseteq P_3$.

\emph{World $1$.} $V+d$ is $V+\tau_{T_{k^*},B}$ with credit at $1$ added on
$Z$, so $V+d\prec_1V+\tau_{T_{k^*},B}\preceq_1V$ by strict (G).

\emph{World $2$.} The change $\delta\coloneqq o_{13}(Q_{k^*})+o_{12}(Z)$ is
debt at $2$ on $Q_{k^*}$ and credit at $2$ on $Z$, inside $D_{23}$, and
lowers $\Lambda_2$ by $\kappa_2P^2_c(Z)-P^2_d(Q_{k^*})>0$; so
$V+\delta\prec_2V$. The pairs $(V+d,V+\tau_{T_{k^*-1},B})$ and $(V+\delta,V)$
differ exactly on $Q_{k^*}\cup Z$ with the same advice there, so by (IB)
$V+d\prec_2V+\tau_{T_{k^*-1},B}\preceq_2V$.

\emph{World $3$.} $o_{13}$ is credit at $3$ on $P_2$, so by (G) and
Lemma~\ref{lem:B1}, $V+d\preceq_3V+o_{12}(Z)+o_{23}(B)\prec_3V$.

So $V+d$ dominates $V$, a contradiction. Hence $V+\tau_{T,B}\not\preceq_1V$
for every $T\subseteq T_0(E)$, and completeness gives the claim.
\end{proof}

\begin{lemma}\label{lem:B3}
Assume $V$ undominated, and let $E$ be a selection on the open edge
$(e_1,e_2)$ with $\beta_E>0$. Then every pair $(A,T)$ of world~$1$ with $A$
non-null and $T\subseteq T_0(E)$ is AC: $T_0(E)$ is credit-null at world~$1$
in the sense of Section~\ref{sec:LH}.
\end{lemma}
\begin{proof}
Suppose $(A,T)$ is SI. By Lemma~\ref{lem:yard}(a) on $[e_1,e_3]$ there are
$A_0\preceq_dA$ and $b$ with $\Band_3(b,\infty)\prec_dA_0$. The set
$Z\coloneqq B_0(E)\cap\Band_3(b,\infty)=P_1\cap\{t_2<\beta_E\}\cap\{t_3>b\}$
is non-null, and $Z\prec_dA$. The context $V+\cred(T)$ is silent on $Z\cup A$
($Z\subseteq P_1$, $T\subseteq P_2$, $A\cap T=\varnothing$), so
$V+\debt(Z)+\cred(T)\prec_1V+\debt(A)+\cred(T)\prec_1V$. On $P_1$, $o_{23}$ is
debt at $1$, and on $P_2$, $o_{13}$ is credit at $1$; so
$V+\tau_{T,Z}\prec_1V$, contradicting Lemma~\ref{lem:B2}.
\end{proof}

\begin{lemma}\label{lem:B4}
Assume $V$ undominated. No selection on the open edge $(e_1,e_2)$ is
two-sided.
\end{lemma}
\begin{proof}
Let $E$ be two-sided. Then $T_0(E)$ and $Q_0\coloneqq P_2\cap\{t_2<\beta_E\}$
are non-null, and $E$ gives $o_2$ (debt at $1$, credit at $2$) on $Q_0$. Put
$p\coloneqq P^2_d(T_0(E))>0$, choose a non-null $Q\subseteq Q_0$ with
$\kappa_2P^2_c(Q)<p$, split $Q$ into non-null $Q',Q''$, and let $X$ be $E$
silenced on $Q\cup T_0(E)$.

\emph{World $1$.} By Lemma~\ref{lem:B3} with $A\coloneqq Q'$ and
$T\coloneqq T_0(E)$, $V\preceq_1V+\debt(Q')+\cred(T_0(E))$, and strict (G)
gives $V\prec_1V+\debt(Q)+\cred(T_0(E))$. The pairs
$(V,V+\debt(Q)+\cred(T_0(E)))$ and $(X,E)$ differ exactly on $Q\cup T_0(E)$
with the same advice, so $X\prec_1E$.

\emph{World $2$.} $E=X+\cred_2(Q)+\debt_2(T_0(E))$, and $X,E$ agree off
$Q\cup T_0(E)\subseteq D_{23}$. So $\Lambda_2(E)-\Lambda_2(X)=
P^2_d(T_0(E))-\kappa_2P^2_c(Q)>0$, and $X\prec_2E$.

So $E$ is beaten at both worlds of its support, contrary to
competitiveness.
\end{proof}

Alternatively, once Theorem~\ref{thm:LH} below is available,
Lemma~\ref{lem:B4} follows from Lemma~\ref{lem:B3} alone: for a two-sided $E$
the set $T_0(E)$ is non-null and credit-null at world~$1$, which
Theorem~\ref{thm:LH} excludes. The proof above uses only the fine edge.

\begin{theorem}\label{thm:fineedge}
Let $n=3$, let all three rankings be complete with (G) and (IB), assume (S),
and suppose $V$ is undominated. If the edge $[e_2,e_3]$ is fine, then
(EP$_{12}$) and (EP$_{13}$) hold.
\end{theorem}
\begin{proof}
A selection on $[e_1,e_2]$ whose credal set omits both endpoints lies in the
open edge and is two-sided; Lemma~\ref{lem:B4} excludes it. (EP$_{13}$)
follows by exchanging the labels $2$ and $3$, under which the hypotheses are
invariant; it uses Lemma~\ref{lem:B1} at world~$2$ and $\Lambda_3$ in
Lemmas~\ref{lem:B2} and~\ref{lem:B4}. So each of (EP$_{12}$), (EP$_{13}$)
uses the local extensions of both worlds $2$ and $3$.
\end{proof}

\begin{corollary}\label{cor:fineedge}
Let $n=3$, let all three rankings be complete with (G) and (IB), and assume
(S). If some edge is fine, $V$ is dominated. Hence, if $V$ is undominated, no
edge is fine and every edge carries a two-sided selection.
\end{corollary}
\begin{proof}
Relabel so that $[e_2,e_3]$ is fine, and suppose $V$ is undominated.
Theorem~\ref{thm:fineedge} gives (EP$_{12}$) and (EP$_{13}$), CCF's
Theorem~7 makes world~$1$ fine, and Theorem~\ref{thm:onefine} makes $V$
dominated: a contradiction. If $V$ is undominated, each edge is non-fine, so
by Lemma~\ref{lem:EPfine} it fails (EP), and a selection omitting both
endpoints is two-sided.
\end{proof}

\begin{remark}\label{rem:edgeroute}
Corollary~\ref{cor:fineedge} does not need Theorem~\ref{thm:onefine}. After
Theorem~\ref{thm:fineedge}, CCF's Theorem~15 gives world~$1$ a loss of
CBOI's form that weakly extends $\preceq_1$ on all regular pairs. By
Lemma~\ref{lem:EPfine}, (EP$_{12}$) makes world~$2$ fine on $D_{12}$; it is
fine on $D_{23}$ by hypothesis, hence fine by Lemma~\ref{lem:fineregion}, and
Corollary~\ref{cor:ext}(b) gives it such a loss. CCF's Theorem~18, with the
complete world~$3$ as the third world, makes $V$ dominated.

Theorem~\ref{thm:fineedge} and CCF's Proposition~17 are independent results.
Proposition~17 assumes only one-sided bounds \eqref{eq:W} at worlds $2$ and
$3$, with neither completeness nor (IB) there. Theorem~\ref{thm:fineedge}
needs worlds $2$ and $3$ complete with (IB) and fine on $D_{23}$, but no
extension beyond $D_{23}$. A loss of CBOI's form does not imply fineness
(Proposition~\ref{ex:germ}). The two results overlap when worlds $2,3$ are
complete with (IB) and fine on $D_{23}$, for instance when they are ranked by
ordinary additive losses as in CCF's Corollary~19.
\end{remark}

The currency of the proof is $o_{12}$ on $P_3$, good at worlds $1$ and $2$
and bad at world~$3$: world~$3$ accepts it against the credit $o_{23}$ gives
it on $P_1$ (Lemma~\ref{lem:B1}), and world~$2$ uses it to absorb the last
piece of debt in Lemma~\ref{lem:B2}. Every world pays on its own odd region.
Without a fine edge, this mechanism leaves exactly the relative fineness
discussed in Theorem~\ref{thm:ownregion}.

\section{One tame world}\label{sec:tame}

\subsection{Pure sets and atoms}

\begin{lemma}\label{lem:YH}
Let $\nu_1,\dots,\nu_r$ be finite nonnegative finitely additive charges on
the measurable subsets of $H$, vanishing on null sets. There are measurable
$Z_1\supseteq Z_2\supseteq\cdots$ with $\mu(Z_m)\to0$ and integrable
$f_l\ge0$ such that $\nu_l(Q)=\int_Qf_l\,d\mu$ for every measurable
$Q\subseteq H\setminus Z_m$, every $m$ and every $l$. If $\nu_l$ is positive
on non-null sets, then $f_l>0$ almost everywhere.
\end{lemma}
\begin{proof}
By the Yosida--Hewitt decomposition \cite{YH} and its concentration
property \cite[Theorem~1 and its proof]{PP}, as used in CBOI's Lemma~3 and
in FAED's equation~(2), $\nu_l=f_l\mu+\pi_l$ with $f_l\ge0$ integrable and
$0\le\pi_l\le\nu_l$, and there are sets $Z'_m$ with $\mu(Z'_m)\to0$ and
$\pi_l(H\setminus Z'_m)=0$ for all $l$ (take unions over $l$). Pass to a
subsequence with $\mu(Z'_m)<2^{-m}$ and put $Z_m\coloneqq\bigcup_{k\ge m}Z'_k$.
Then $Z_m$ decreases, $\mu(Z_m)\le2^{1-m}$, and $\pi_l$ vanishes on subsets
of $H\setminus Z_m\subseteq H\setminus Z'_m$. If $f_l=0$ on a non-null $N$,
choose $m$ with $\mu(Z_m)<\mu(N)$; then $\nu_l(N\setminus Z_m)=0$ with
$N\setminus Z_m$ non-null.
\end{proof}

We call such $Z_m$ \emph{pure sets} for $\nu_1,\dots,\nu_r$. Charges defined
only on the subsets of a measurable set are extended by $0$. For a charge
$\nu$ and an ultrafilter $\Uf$ put $\nu(\Uf)\coloneqq\inf\{\nu(B):B\in\Uf\}$,
as in Definition~\ref{def:partner}.

\begin{lemma}\label{lem:atoms}
Let $\nu$ have pure sets $Z_m$. If an ultrafilter $\Uf$ contains a set $R$
with $R\cap Z_m$ null for some $m$, then $\nu(\Uf)=0$.
\end{lemma}
\begin{proof}
By Lemma~\ref{lem:uf}(a) there are $B_k\in\Uf$ with $\mu(B_k)\to0$. Then
$B_k\cap R\in\Uf$ and $\nu(B_k\cap R)=\int_{B_k\cap R}f\,d\mu\to0$.
\end{proof}

\subsection{Fine off the partner's pure sets}

In this subsection the standing hypotheses of Section~\ref{sec:partner}
hold, world $j$ has a partner bound $(a,b)$ on $D=D_{1j}$, $Z_m$ are pure
sets for $a,b$, and $R_m\coloneqq D\setminus Z_m$.

\begin{corollary}\label{cor:fineoff}
For every $m$ and all non-null $A,T\subseteq H'$, $R_m\in\Jid_A\cap\Kid_T$.
In particular, if $R_m$ is non-null, world~$1$ is fine on $R_m$, and the
almost agreeing probabilities $P_d,P_c$ of its debt and credit orders on
$R_m$ exist and satisfy Proposition~\ref{prop:localext}(i)--(iii) on $R_m$.
\end{corollary}
\begin{proof}
If $R_m\notin\Jid_A$, Lemma~\ref{lem:uf}(b) gives an ultrafilter $\Uf$
containing $R_m\subseteq D$ and avoiding $\Jid_A$; then $a(\Uf)>0$ by
Theorem~\ref{thm:transfer}(a), while $a(\Uf)=0$ by Lemma~\ref{lem:atoms}.
For $\Kid_T$ use Theorem~\ref{thm:transfer}(b) and Lemma~\ref{lem:atoms} for
$a$ and $b$. The last statement uses only fineness: see the proof of
Proposition~\ref{prop:localext}.
\end{proof}

\begin{lemma}\label{lem:pairsoff}
For all large $m$ there are an SI pair and an AC pair of world~$1$,
each with non-null debt and credit parts contained in $R_m$.
\end{lemma}
\begin{proof}
Since $R_m$ increases with $m$, it suffices to find each pair inside some
$R_m$.

\emph{SI pair.} If (EP$_{1j}$) holds, fix $m_0$ and a non-null
$T_0\subseteq\Band_j(1,2)\setminus Z_{m_0}$. The proof of the (EP) case of
Lemma~\ref{lem:SIpartner}, with $T_0$ in place of $T_y$, gives $s<1$ with
$(\Band_j[s,1),T_0)$ SI. Choose $m\ge m_0$ with
$\mu(Z_m)<\mu(\Band_j[s,1))$; then $(\Band_j[s,1)\setminus Z_m,T_0)$ is
SI by (G), and both parts lie in $R_m$. If (EP$_{1j}$) fails, fix a
two-sided selection $C_0$ and $m$ with $\cred(C_0)\cap R_m$ and
$\debt(C_0)\cap R_m$ non-null, and run the non-(EP) case of
Lemma~\ref{lem:SIpartner} with a non-null $T_0\subseteq\cred(C_0)\cap R_m$
in place of $T_y$ and $D''$ chosen in $\debt(C_0)\cap R_m$; that argument
does not use that $T_y$ is a yardstick.

\emph{AC pair.} Let $q_n\uparrow q_e$ and $C_n=(\alpha_n,\beta_n)$ be as in
Lemma~\ref{lem:ends}(b). If $\alpha_n<\infty$ for infinitely many $n$,
choose $m_0$ and a non-null $A''\subseteq\Band_j(0,1)\setminus Z_{m_0}$; by
Lemma~\ref{lem:SAC} there is $n$ with $(A'',W)$ strictly AC for every
non-null $W\subseteq\Band_j(\alpha_n,\infty)$; take $m\ge m_0$ with
$W\coloneqq\Band_j(\alpha_n,\infty)\setminus Z_m$ non-null. Otherwise
$C\coloneqq C_n=E_j(\beta_n)$ for some $n$; take $m$ with
$\Band_j(\beta_n,\infty)\setminus Z_m$ non-null, a non-null $A$ in it
whose complement there is non-null, and
(FAED's Lemma~1) a non-null $T\subseteq\Band_j(\beta_n,\infty)\setminus(Z_m
\cup A)$ with $b(T)<a(A)$. Let $X\coloneqq C+\debt(A)+\cred(T)$; $C$ is silent
on $A\cup T$. By \eqref{eq:incr}, $\Lambda_j(X)-\Lambda_j(C)=b(T)-a(A)<0$, so
$X\prec_jC$, $C\preceq_1X$, and by (IB) $(A,T)$ is AC.
\end{proof}

\subsection{The chain lemma}

For the rest of this section $n=3$ and $\{u,v,w\}=\{1,2,3\}$. On $P_w$ the
option $o_{uv}$ is good at $u$ and $v$; on $P_v$, $o_{uw}$ is credit at $u$ and
$w$ and debt at $v$; on $P_u$, $o_{vw}$ is credit at $v$ and $w$ and debt at
$u$. For a selection $C$ on $[e_u,e_v]$ put
\[
 T_0(C)\coloneqq P_v\cap\{C\text{ gives }o_{uw}\},\qquad
 B_0(C)\coloneqq P_u\cap\{C\text{ gives }o_{vw}\}.
\]
For $(u,v)=(1,2)$ these are the sets of Section~\ref{sec:edge}. If $C$ is
two-sided, both are non-null: in world $u$'s frame $T_0(C)=P_v\cap
\{t^u_v>\alpha_C\}$ and $B_0(C)=P_u\cap\{t^u_v<\beta_C\}$.

\begin{lemma}\label{lem:chain}
Assume (IB) at worlds $u$ and $v$ and (G) at world~$u$. Let $V$ be
competitive on $\{1,2,3\}$ and let $C$ be competitive on $\{u,v\}$. Let
$A\subseteq D_{uv}$ and $S\subseteq P_w$ be measurable with $S$ non-null,
$X'\coloneqq V+o_{uv}(S)$, $X$ the model $X'$ with $C$'s advice on $A$, and $Y$
the model $C$ silenced on $A$. If $V\not\preceq_vX$, $V\not\preceq_wX$ and
$X\not\preceq_vX'$, then strict (G) fails.
\end{lemma}
\begin{proof}
This is IC's Lemma~2 with the labels $(1,2,3)$ replaced by $(u,v,w)$; its
proof is order-theoretic. By competitiveness of $V$, $V\preceq_uX$. The
pairs $(C,Y)$ and $(X,X')$ differ exactly on $A$ with the same advice there
($X'$ is silent on $A$, since $S\cap A=\varnothing$). By (IB) at $v$,
$X\not\preceq_vX'$ gives $C\not\preceq_vY$; competitiveness of $C$ gives
$C\preceq_uY$; and (IB) at $u$ gives $X\preceq_uX'$. So
$V\preceq_uX\preceq_uX'$. But $X'\prec_uV$ by strict (G), since $o_{uv}$ is
good at $u$ on the non-null $S$.
\end{proof}

\subsection{One world with extensions on both of its edges}

\emph{Standing for this subsection:} $n=3$; all three rankings are complete
with (G) and (IB); (S); and $V$ is undominated.

\begin{theorem}\label{thm:chainEP}
Suppose world $v$ has strictly positive two-sided local extensions of
CBOI's form on both $D_{uv}$ and $D_{vw}$: charges $\nu^{uv}_1,\nu^{uv}_2$ on
$D_{uv}$ and $\nu^{vw}_1,\nu^{vw}_2$ on $D_{vw}$, vanishing on null sets
and positive on non-null sets,
such that for regular $X,Y$ agreeing off $D_{uv}$, $X\preceq_vY$ implies
$\Lambda_{uv}(X)\le\Lambda_{uv}(Y)$, where
$\Lambda_{uv}(M)\coloneqq\nu^{uv}_1(D_{uv}\cap\{r_v\ge1\})
+\nu^{uv}_2(D_{uv}\cap\{r_v\ge2\})$, and likewise for $D_{vw}$. Then
(EP$_{uv}$) holds.
\end{theorem}
\begin{proof}
Suppose a selection $C$ on $[e_u,e_v]$ omits both endpoints; it is two-sided,
in the open edge, and competitive on $\{u,v\}$, and $V$ is competitive on
$\{1,2,3\}$. Write $T_0\coloneqq T_0(C)$ and $B_0\coloneqq B_0(C)$. By the
table above, changes from silence change $\Lambda_{uv}$ and $\Lambda_{vw}$ as
follows: $o_{vw}$ on $E\subseteq P_u$ adds $-\nu^{uv}_1(E)$; $o_{uw}$ on
$E\subseteq P_v$ adds $+\nu^{uv}_2(E)$ to $\Lambda_{uv}$ and $+\nu^{vw}_2(E)$ to
$\Lambda_{vw}$; $o_{uv}$ on $E\subseteq P_w$ adds $-\nu^{vw}_1(E)$. By
completeness, each local extension has the strict form: $\Lambda(X)<\Lambda(Y)$
implies $X\prec_vY$ for pairs agreeing off its region.

Let $Z_m$ be common pure sets for the four charges (Lemma~\ref{lem:YH}), fix
$m_1$ with $\mu(Z_{m_1})<\min(\mu(T_0),\mu(B_0))/2$, put $Z\coloneqq Z_{m_1}$,
$T'_0\coloneqq T_0\setminus Z$ and $B'_0\coloneqq B_0\setminus Z$.

\emph{Step 1.} Choose a non-null $T\subseteq T'_0$ with
$\nu^{uv}_2(T)<\nu^{uv}_1(B'_0)$ and $\mu(T)<\mu(T'_0)/2$; this is possible
because $\nu^{uv}_2$ has a density on $T'_0$.

\emph{Step 2 (the third world).} We claim: there is a non-null
$S\subseteq P_w$ with $V+o_{uv}(S)+o_{uw}(T)\prec_wV$. On $P_w$, $o_{uv}$ is
debt at $w$; on $P_v$, $o_{uw}$ is credit at $w$; both sets lie in $D_{vw}$.
Apply Section~\ref{sec:transfer} and Corollary~\ref{cor:fineoff} with base
world~$w$ and partner $v$ on $D_{vw}$: world $v$'s extension on $D_{vw}$ is a
partner bound by Lemma~\ref{lem:sources}(c), with $a\coloneqq\nu^{vw}_1$ and
$b\coloneqq\nu^{vw}_2$, since by \eqref{eq:swap}
$D_{vw}\cap\{r_v\ge1\}=D_{vw}\setminus\debt_w$ and
$D_{vw}\cap\{r_v\ge2\}=D_{vw}\cap\cred_w$. By Lemma~\ref{lem:pairsoff} there
are $m\ge m_1$ with $\mu(Z_m)<\mu(P_w)/2$ and an SI pair $(A_0,T'')$ of
world~$w$ with non-null parts in $R\coloneqq D_{vw}\setminus Z_m$; by
Corollary~\ref{cor:fineoff} world~$w$ is fine on $R$, with almost agreeing
probabilities $P_d,P_c$ on $R$. Since $Z_m\subseteq Z_{m_1}$,
$T\subseteq P_v\setminus Z_{m_1}\subseteq R$. If $A_0\setminus T$ is
non-null, replace $A_0$ by it (less debt keeps the pair SI, by (G));
otherwise $A_0\subseteq T$ up to a null set, and we replace $A_0$ by a
non-null half $A_1$ of $A_0$ and $T$ by $T\setminus A_1$, which is non-null
and still satisfies Step~1. Now $A_0\cap T=\varnothing$. Choose an integer
$K\ge1$ and $\gamma>0$ with $P_c(T'')/K+\gamma<P_c(T)$ (note $P_c(T)>0$).
Split $A_0$ into $K$ disjoint non-null pieces $A_l$ and, using
Proposition~\ref{prop:localext}(iii), $T''$ into $K$ disjoint pieces $T_l$
with $P_c(T_l)\le P_c(T'')/K+\gamma$. If all pairs $(A_l,T_l)$ were AC, so
would be their composition $(A_0,T'')$; so some $(A_l,T_l)$ is SI. Since
$P_c(T_l)<P_c(T)$, $T_l\prec^w_cT$, and in the context
$V+\debt_w(A_l)$, silent on $T\cup T_l$,
\[
 V+\debt_w(A_l)+\cred_w(T)\prec_wV+\debt_w(A_l)+\cred_w(T_l)\prec_wV.
\]
Choose a non-null $S\subseteq P_w\cap R$ with $0<P_d(S)<P_d(A_l)$; then
$S\prec^w_dA_l$, and in the context $V+\cred_w(T)$, silent on $S\cup A_l$,
$V+\debt_w(S)+\cred_w(T)\prec_wV+\debt_w(A_l)+\cred_w(T)\prec_wV$. This is
the claim.

\emph{Step 3.} Choose a non-null $Q\subseteq T'_0\setminus T$ with
$\nu^{vw}_2(Q)<\nu^{vw}_1(S)$; $\nu^{vw}_2$ has a density on $P_v\setminus Z$,
and $T'_0\setminus T$ is non-null by Step~1. The models
$V+o_{uw}(Q)+o_{uv}(S)$ and $V$ agree off $Q\cup S\subseteq D_{vw}$, and
$\Lambda_{vw}$ changes by $\nu^{vw}_2(Q)-\nu^{vw}_1(S)<0$, so
$V+o_{uw}(Q)+o_{uv}(S)\prec_vV$.

\emph{Step 4.} Let $\Sigma_s$, $s\in[0,1]$, be an increasing family of
subsets of $B'_0$ with $\Sigma_0=\varnothing$, $\Sigma_1=B'_0$ and
$\mu(\Sigma_s\triangle\Sigma_r)\to0$ as $r\to s$ (for instance
$\Sigma_s\coloneqq B'_0\cap\{t^u_v<\tan(\pi s/2)\}$). Put
$\tau_s\coloneqq o_{uw}(T)+o_{vw}(\Sigma_s)$ and
\[
 \lambda(s)\coloneqq\Lambda_{uv}(V+\tau_s)-\Lambda_{uv}(V)
 =\nu^{uv}_2(T)-\nu^{uv}_1(\Sigma_s).
\]
$\lambda$ is continuous, since $\nu^{uv}_1$ has a density on $B'_0$;
$\lambda(0)>0$, and $\lambda(1)<0$ by Step~1. Choose $s$ with
$0<\lambda(s)<\nu^{uv}_2(Q)$, and put $X'\coloneqq V+o_{uv}(S)$ and
$X\coloneqq X'+\tau_s$.
\begin{itemize}[leftmargin=*,itemsep=1pt]
\item $X'\prec_vX$: the two agree off $T\cup\Sigma_s\subseteq D_{uv}$, $X'$
agrees with $V$ on $D_{uv}$, and $\Lambda_{uv}(X)-\Lambda_{uv}(X')=
\lambda(s)>0$. So $X\not\preceq_vX'$.
\item $X\prec_vV$: $V+\tau_s$ and $V+o_{uw}(Q)$ agree off
$T\cup\Sigma_s\cup Q\subseteq D_{uv}$, with $\Lambda_{uv}$-difference
$\lambda(s)-\nu^{uv}_2(Q)<0$, so $V+\tau_s\prec_vV+o_{uw}(Q)$. The pairs
$(V+\tau_s,V+o_{uw}(Q))$ and $(X,V+o_{uw}(Q)+o_{uv}(S))$ differ exactly on
the pairwise disjoint sets $T,\Sigma_s,Q$ with the same advice, so by (IB)
and Step~3, $X\prec_vV+o_{uw}(Q)+o_{uv}(S)\prec_vV$. So $V\not\preceq_vX$.
\item $X\prec_wV$: $X$ is $V+o_{uv}(S)+o_{uw}(T)$ with $o_{vw}$, good at $w$,
added on $\Sigma_s\subseteq P_u$; by (G) and Step~2,
$X\preceq_wV+o_{uv}(S)+o_{uw}(T)\prec_wV$. So $V\not\preceq_wX$.
\end{itemize}

\emph{Step 5.} $X$ is $X'$ with $C$'s advice on $A\coloneqq T\cup\Sigma_s$:
$C$ gives $o_{uw}$ on $T\subseteq T_0$ and $o_{vw}$ on $\Sigma_s\subseteq B_0$.
Lemma~\ref{lem:chain} gives a contradiction.
\end{proof}

The strict form of both local extensions is used at pairs of models that are
not selections (Steps~3 and~4), so the theorem needs two-sided extensions.
World $u$ enters only through Lemma~\ref{lem:chain} and the yardsticks, so it
needs only (G) and (IB). The next theorem uses a \emph{one-sided} bound
instead, at the price of making it global.

\begin{theorem}\label{thm:chainW}
Suppose world $v$ has a loss $L_v$ of CBOI's form with \eqref{eq:positive}
and \eqref{eq:W}. Then (EP$_{uv}$) holds.
\end{theorem}
\begin{proof}
Let $\nu_1,\nu_2$ be $L_v$'s charges on $H'$. By Lemma~\ref{lem:sources}(b)
their restrictions to $D_{vw}$ form a partner bound of world $v$ for the base
world $w$, so Corollary~\ref{cor:fineoff} and Lemma~\ref{lem:pairsoff} apply
with base~$w$. Suppose a selection $C$ on $[e_u,e_v]$ omits both endpoints,
and carry out Steps~1 and~2 of the proof of Theorem~\ref{thm:chainEP} with
$\nu^{uv}_k\coloneqq\nu_k|_{D_{uv}}$ and pure sets for $\nu_1,\nu_2$. This gives
non-null $T\subseteq T_0\setminus Z$ with $\nu_2(T)<\nu_1(B_0\setminus Z)$ and
non-null $S\subseteq P_w$ with $V+o_{uv}(S)+o_{uw}(T)\prec_wV$. Define
$\Sigma_s$ and $\tau_s$ as in Step~4, and $\lambda(s)\coloneqq\nu_2(T)
-\nu_1(\Sigma_s)$; choose $s$ with $0<\lambda(s)<\nu_1(S)$. Put
$X'\coloneqq V+o_{uv}(S)$, $X\coloneqq X'+\tau_s$, and let $Y$ be $C$
silenced on $T\cup\Sigma_s$. At world $v$, $o_{uv}$ on $S$ and $o_{vw}$ on
$\Sigma_s$ are credit and $o_{uw}$ on $T$ is debt, so
\[
 L_v(X)-L_v(V)=\lambda(s)-\nu_1(S)<0,\qquad L_v(Y)-L_v(C)=-\lambda(s)<0.
\]
By \eqref{eq:W} at the coherent undominated $V$ and at the selection $C$,
$V\not\preceq_vX$ and $C\not\preceq_vY$; by (IB) the latter is
$X\not\preceq_vX'$. As in Step~4, $V\not\preceq_wX$. Lemma~\ref{lem:chain}
gives a contradiction.
\end{proof}

The world-$v$ computation is that of CCF's Proposition~17; the new
ingredient is the third world's exchange in Step~2, which comes from world
$w$'s fineness off $v$'s pure sets. A merely local one-sided bound would not
suffice: $V$ and $X$ differ on $D_{uv}\cup P_w$.

\subsection{Tame worlds}

Let $(\Gamma,+,\le)$ be an ordered abelian group. A \emph{group primary} at
world $v$ is a function $L(X)\coloneqq\nu_1(H'\cap\{r_v(\cdot,X)\ge1\})+
\nu_2(H'\cap\{r_v(\cdot,X)\ge2\})$ on regular models, with
$\nu_1,\nu_2\ge0$ finitely additive $\Gamma$-valued maps on the measurable
subsets of $H'$ vanishing on null sets, such that $L(X)<L(Y)$ implies
$X\prec_vY$. A \emph{local} group primary on $R$ is the same with $H'$
replaced by $R$ and $X,Y$ restricted to pairs agreeing off $R$. Examples are
exact representations by lexicographic hierarchies of signed finitely
additive criteria, by hyperreal-valued charges, or by ultrapowers of
measures. Let $W_v$ be the regular model that gives debt at $v$ on all of
$H'$.

\begin{lemma}\label{lem:leading}
If world $v$ has a group primary $L$ that is nonconstant on regular models,
it has a real-valued primary $\Phi$ of FAED's form \eqref{eq:primaryFAED}
below, with $\Phi(W_v)-\Phi(D_v)=1$. If (G) and (S) hold, $\Phi$'s charges
are positive on non-null sets, and world $v$ has a loss of CBOI's form with
\eqref{eq:positive} that weakly extends $\preceq_v$ on all regular pairs.
The same holds locally: a nonconstant local group primary on $D_{vw}$ gives a
strictly positive two-sided local extension of CBOI's form on $D_{vw}$.
\end{lemma}
Here FAED's form, from FAED's equation (1), is
\begin{equation}
 \Phi(X)=\nu'_1(H'\cap\{r_v(\cdot,X)\ge1\})+\nu'_2(H'\cap\{r_v(\cdot,X)\ge2\})
 \label{eq:primaryFAED}
\end{equation}
with finite nonnegative real finitely additive $\nu'_k$ vanishing on null
sets, and a \emph{primary} means $\Phi(X)<\Phi(Y)\Rightarrow X\prec_vY$.
\begin{proof}
$L(D_v)=0$, and for regular $X$, $0\le L(X)\le M\coloneqq L(W_v)=\nu_1(H')+
\nu_2(H')$ by finite additivity and $\nu_k\ge0$. Since $L$ is nonconstant,
$M>0$. Convex subgroups of $\Gamma$ are totally ordered by inclusion; let
$C^+$ be the smallest one containing $M$ and $C_-$ the largest one not
containing $M$. Every convex subgroup between them is one of them, so
$C^+/C_-$ has no nontrivial proper convex subgroup and is Archimedean. By
H\"older's theorem \cite{Holder} there is an order-preserving homomorphism
$h\colon C^+/C_-\to\R$ with $h(M+C_-)=1$. By convexity $[0,M]\subseteq C^+$.
Put $\varphi(x)\coloneqq h(x+C_-)$ for $x\in C^+$; $\varphi$ is additive and
monotone. Then $\nu'_k\coloneqq\varphi\circ\nu_k$ are real, nonnegative,
finitely additive, bounded by $1$ and null on null sets, and
$\Phi\coloneqq\varphi\circ L$ has the form \eqref{eq:primaryFAED}. If
$\Phi(X)<\Phi(Y)$, then $L(X)<L(Y)$ (otherwise $L(Y)\le L(X)$ would give
$\Phi(Y)\le\Phi(X)$), so $X\prec_vY$. And $\Phi(W_v)-\Phi(D_v)=1$.

Under (G) and (S), FAED's Theorem~3 applies to the nonconstant primary
$\Phi$ and gives alternative (a): the ordinary parts of $\nu'_1,\nu'_2$ have
densities that are positive almost everywhere on $H'$, so $\nu'_k$ is
positive on non-null sets. A primary weakly extends $\preceq_v$: if
$X\preceq_vY$ then not $Y\prec_vX$, so $\Phi(X)\le\Phi(Y)$. Put
$\nu_{v1}\coloneqq\nu'_1+\mu|_{H\setminus H'}$,
$\nu_{v2}\coloneqq\nu'_2+\mu|_{H\setminus H'}$ and $\nu_{v3}\coloneqq\mu$;
on regular models the resulting loss is $\Phi$ plus a constant.

Locally, the same construction with $H'$ replaced by $D_{vw}$ and $W_v$ by
the model with debt at $v$ on $D_{vw}$ gives a real primary for pairs
agreeing off $D_{vw}$. FAED's Lemma~2 and the part of FAED's Theorem~3 that
concerns the single edge $[e_v,e_w]$ compare only edge-$vw$ selections with
regular models that differ from them on subsets of $D_{vw}$; all these agree
off $D_{vw}$. So they apply, and the blind alternative is excluded because
$\nu'_1(D_{vw})+\nu'_2(D_{vw})=1$.
\end{proof}

World $v$ satisfies the \emph{finite cancellation conditions} (FC) if there
are no regular models $X_1,\dots,X_m$ and $Y_1,\dots,Y_m$, repetitions
allowed, with
$X_l\preceq_vY_l$ for all $l$, $X_l\prec_vY_l$ for some $l$, and
$\sum_l(u^{X_l}-u^{Y_l})\ge0$ almost everywhere on $H'\times\{1,2\}$, where
\[
 u^X\coloneqq\bigl(\mathbf 1\{r_v(\cdot,X)\ge1\},\ \mathbf 1\{r_v(\cdot,X)\ge2\}\bigr).
\]
The local form on $R$ restricts to pairs agreeing off $R$ and to
$R\times\{1,2\}$.

\begin{lemma}\label{lem:FC}
If world $v$ is complete and satisfies (FC), it has a group primary $L$ with
$X\preceq_vY\iff L(X)\le L(Y)$ for all regular $X,Y$. If (G) holds at $v$,
$L$ is nonconstant. The same holds in the local form.
\end{lemma}
\begin{proof}
Consider the first-order language of ordered abelian groups with a constant
$c_{k,B}$ for each $k\in\{1,2\}$ and each element $B$ of the measure algebra
$\mathfrak A$ of $H'$. Let $\Sigma$ consist of the axioms of ordered abelian
groups, $c_{k,B}\ge0$, $c_{k,\varnothing}=0$, $c_{k,B\cup C}=c_{k,B}+c_{k,C}$
for disjoint $B,C$, and, writing
$L(X)\coloneqq c_{1,B_1(X)}+c_{2,B_2(X)}$ with
$B_k(X)\coloneqq H'\cap\{r_v(\cdot,X)\ge k\}$, the
sentences $L(X)\le L(Y)$ for $X\preceq_vY$ and $L(X)<L(Y)$ for $X\prec_vY$.
A model of $\Sigma$ gives $\Gamma$ and $\nu_k(B)\coloneqq c_{k,B}$ as required;
by completeness, $X\preceq_vY$ or $Y\prec_vX$, so $L$ represents
$\preceq_v$ exactly. By the compactness theorem it suffices to satisfy each
finite $\Sigma_0\subseteq\Sigma$ in $(\R,+,\le)$. Let $S_1,\dots,S_p$ be the
non-null atoms of the finite subalgebra generated by the sets occurring in
$\Sigma_0$, and interpret $c_{k,B}\coloneqq\sum_{S_i\subseteq B}x_{k,i}$ with
unknowns $x_{k,i}\ge0$. The axioms of additivity, emptiness and positivity
then hold, and the comparisons in $\Sigma_0$ become finitely many weak
inequalities $(u^{X_l}-u^{Y_l})\cdot x\le0$ and strict inequalities
$(u^{X_l}-u^{Y_l})\cdot x<0$ with integer coefficients ($u^X$ is constant on
each $S_i\times\{k\}$). If there is no strict one, $x=0$ works. Otherwise,
by Motzkin's transposition theorem \cite{Motzkin}, the system with $x\ge0$
is infeasible only if a nonnegative combination of its rows, with positive weight on some
strict row, is $\ge0$ coordinatewise; the data are rational, so the weights
may be taken to be integers. Repeating each pair according to its weight
then gives a system that violates (FC). So $\Sigma_0$ is satisfiable. If (G)
holds, $D_v\prec_vW_v$, so $L(D_v)<L(W_v)$. The local form is the same with
pairs agreeing off $R$ and $H'$ replaced by~$R$.
\end{proof}

Conversely an exact group-valued representation satisfies (FC): summing the
inequalities $L(X_l)\le L(Y_l)$, one of them strict, contradicts
$\sum_l(u^{X_l}-u^{Y_l})\ge0$ and $\nu_k\ge0$.

\begin{theorem}\label{thm:tame}
Let $n=3$, let all three rankings be complete with (G) and (IB), and assume
(S). Suppose some world $v$ satisfies one of the following:
\begin{enumerate}[label=(\roman*),leftmargin=*,itemsep=1pt]
\item world $v$ is fine;
\item world $v$ has a loss of CBOI's form with \eqref{eq:positive} and
\eqref{eq:W}, that is, by EDL's Theorem~2, world $v$ carries no leverage
sequence; this includes the case that world $v$ has a strictly
positive finitely additive weak extension of CBOI's form on all regular
pairs;
\item world $v$ has a nonconstant group primary;
\item world $v$ satisfies (FC);
\item world $v$ has strictly positive two-sided local extensions of CBOI's
form on both of its edge regions; this holds in particular if world
$v$ has nonconstant local group primaries on both, or satisfies the local
(FC) on both.
\end{enumerate}
Then $V$ is dominated, by a selected coherent model.
\end{theorem}
\begin{proof}
Suppose $V$ is undominated. (i) contradicts Theorem~\ref{thm:onefine}. In
case (ii), Theorem~\ref{thm:chainW}, applied with the roles of $u$ and $w$ in
both orders, gives (EP$_{uv}$) and (EP$_{wv}$), so world $v$ is fine by CCF's
Theorem~7, which is case (i). (iii) implies (ii) by Lemma~\ref{lem:leading},
and (iv) implies (iii) by Lemma~\ref{lem:FC}. In case (v),
Theorem~\ref{thm:chainEP} gives (EP) on both edges at $v$, and again world
$v$ is fine; the two special cases give local extensions by
Lemmas~\ref{lem:FC} and~\ref{lem:leading}. So $V$ is dominated; by CBOI's
Lemma~1, by a selected coherent model.
\end{proof}

\begin{corollary}\label{cor:tameshape}
Let $n=3$, let all three rankings be complete with (G) and (IB), assume (S),
and suppose $V$ is undominated. Then every world carries a leverage sequence
in EDL's sense, has no nonconstant group primary, violates some finite
cancellation condition, and is not fine. For every world $v$, at least one of
its edge regions carries no strictly positive two-sided local extension of
world~$v$. If world $v$ has such an extension on $D_{vx}$, then world $x$ is
fine on $D_{vx}\setminus Z_m$ for the pure sets $Z_m$ of $v$'s charges, and
every ultrafilter on $D_{vx}$ that avoids some $\Jid^x_A$ or $\Kid^x_T$ is an
atom of $v$'s charges.
\end{corollary}
\begin{proof}
Theorem~\ref{thm:tame}, Corollary~\ref{cor:fineoff} and
Theorem~\ref{thm:transfer} with base~$x$ and partner~$v$
(Lemma~\ref{lem:sources}(c)).
\end{proof}

\noindent This strengthens CCF's Corollary~20(a), by which at most one world
has a loss with \eqref{eq:W}, to: no world has one.

\subsection{Extensions off small sets}

\begin{definition}\label{def:aext}
World $x$ has \emph{(AExt) on $D_{xy}$} if there are measurable
$Z_1\supseteq Z_2\supseteq\cdots$ with $\mu(Z_m)\to0$ and, for each $m$,
charges $\nu^m_1,\nu^m_2$ on the measurable subsets of
$R_m\coloneqq D_{xy}\setminus Z_m$, vanishing on null sets and
positive on non-null sets, such that for
regular $X,Y$ agreeing off $R_m$, $X\preceq_xY$ implies
$\Lambda^m(X)\le\Lambda^m(Y)$, where $\Lambda^m(M)\coloneqq
\nu^m_1(R_m\cap\{r_x\ge1\})+\nu^m_2(R_m\cap\{r_x\ge2\})$.
\end{definition}

So world $x$ is Archimedean, in the sense of having a finitely additive
extension, once a set of arbitrarily small measure is removed. A
two-sided local extension on $D_{xy}$ gives (AExt), with $Z_m$ its pure sets.
We may assume that each $\nu^m_k$ has a density on $R_m$ (``with
densities''): pass to a subsequence with $\sum_m\mu(Z_m)<\infty$, let $Y_m$
be a pure set for $\nu^m_1,\nu^m_2$ with $\mu(Y_m)<2^{-m}$, replace $Z_m$ by
$Z'_m\coloneqq\bigcup_{k\ge m}(Z_k\cup Y_k)$, and restrict $\nu^m_k$ to
$D_{xy}\setminus Z'_m$. Pairs agreeing off $D_{xy}\setminus Z'_m$ agree off
the old $R_m$, and their $\Lambda^m$-difference involves only subsets of
$D_{xy}\setminus Z'_m$. If world $x$ is complete, $\Lambda^m$ has the strict
form.

\begin{theorem}\label{thm:aexttransfer}
Assume the standing hypotheses of Section~\ref{sec:partner}, and let world
$j$ have (AExt) on $D=D_{1j}$, with densities. Then:
\begin{enumerate}[label=(\alph*)]
\item $R_m\in\Jid_A\cap\Kid_T$ for every $m$ and all non-null $A,T\subseteq
H'$;
\item for all large $m$ there are an SI pair and an AC pair of
world~$1$ with non-null debt and credit parts contained in $R_m$.
\end{enumerate}
\end{theorem}
\begin{proof}
In this proof the index of (AExt) is written $r$: the data are $Z_r$,
$R_r$, $\nu^r_1,\nu^r_2$ and $\Lambda^r$. The indices $m$, $m'$, $m_0$ keep
their meaning from the proof of Theorem~\ref{thm:transfer} (the selections
$C_m$ and the threshold of Lemma~\ref{lem:SAC}), and the yardstick $Z$ of
Steps~1(ii) and~2(ii) of that proof is called $Y$ here.

(a) Suppose an ultrafilter $\Uf$ contains $R_r$ and avoids $\Jid_A$
(resp.\ $\Kid_T$). Then $\Uf\ni R_{r'}$ for all $r'\ge r$, and for
every $R\in\Uf$ and $\varepsilon>0$ some $B\in\Uf$, $B\subseteq R\cap
R_{r'}$, has $\nu^{r'}_1(B)+\nu^{r'}_2(B)<\varepsilon$, by
Lemma~\ref{lem:uf}(a) and the densities. Run the proof of
Theorem~\ref{thm:transfer}(a) (resp.\ (b)), including
Lemmas~\ref{lem:creditconf} and~\ref{lem:SIpartner}, with $a,b$ replaced at
each use by the charges of a suitable $\Lambda^{r'}$, as follows. Each
use of the partner bound there compares a selection $C$ (one of $C_m$,
$C_0$, $E$) with a regular model that differs from $C$ exactly on one piece
chosen in $\Uf$ ($B$, or $W$ in Step~2(ii)), or on the set $D''$ of
Lemma~\ref{lem:SIpartner}, and on one auxiliary set fixed before it ($B'$,
$W$, $Y$ or $T_y$). (1) At the start replace $T_y$ by
$T_y\setminus Z_{r_0}$, where $r_0\ge r$ and
$\mu(Z_{r_0})<\mu(T_y)/2$; it is still a non-null credit yardstick
$\preceq_cT$ inside $D\cap\{t\in J\}$, and this is done before $D'$ is
constructed. (2) At each use, once the auxiliary set is fixed, choose
$r'\ge r_0$ with $\mu(Z_{r'})$ less than half its measure, and replace
it by its part off $Z_{r'}$: $B'\setminus Z_{r'}$ is still a yardstick
in the same region, $W\setminus Z_{r'}$ keeps the pair $(A',W)$
strictly AC by (G), and $Y\setminus Z_{r'}$ is still $\prec_dD'$. (3)
Choose the piece (or $D''$) inside $R_{r'}$ with small
$\Lambda^{r'}$-charges. The two compared models then agree off
$R_{r'}$, so $\Lambda^{r'}$ and its strict form apply, and every
sign computation is unchanged.

(b) The proof of Lemma~\ref{lem:pairsoff} with all sets chosen inside
$R_r$ for large $r$ and with $\Lambda^r$ in place of $\Lambda_j$.
\end{proof}

\begin{theorem}\label{thm:aext}
Let $n=3$, let all three rankings be complete with (G) and (IB), and assume
(S). If some world $v$ has (AExt) on both of its edge regions, $V$ is
dominated.
\end{theorem}
\begin{proof}
Take $Z_m\coloneqq Z^{uv}_m\cup Z^{vw}_m$, with densities; both families of
extensions still apply off $Z_m$. Suppose $V$ is undominated, and repeat the
proof of Theorem~\ref{thm:chainEP} with these changes. In Step~1 use
$\Lambda^{uv,m_1}$ on $T\cup B'_0\subseteq D_{uv}\setminus Z_{m_1}$. In
Step~2 use Theorem~\ref{thm:aexttransfer} with base $w$ and partner $v$ on
$D_{vw}$: world $w$ is fine on $R\coloneqq D_{vw}\setminus Z_m$ and has an SI
pair inside $R$ for some $m\ge m_1$ with $\mu(Z_m)<\mu(P_w)/2$; the scaling
argument is unchanged and gives $S\subseteq P_w\cap R$. In Step~3 choose
$Q\subseteq T'_0\setminus T$ with $\nu^{vw,m}_2(Q)<\nu^{vw,m}_1(S)$; the
compared models agree off $Q\cup S\subseteq D_{vw}\setminus Z_m$, since
$Z_m\subseteq Z_{m_1}$. In Step~4 use $\Lambda^{uv,m_1}$; all models compared
there agree off $T\cup\Sigma_s\cup Q\subseteq D_{uv}\setminus Z_{m_1}$. This
gives (EP$_{uv}$), symmetrically (EP$_{wv}$), and world $v$ is fine by CCF's
Theorem~7, so Theorem~\ref{thm:onefine} gives a contradiction.
\end{proof}

\begin{theorem}\label{thm:aextshared}
Assume the standing hypotheses of Section~\ref{sec:partner}. If world $j$ has
(AExt) on $D_{1j}$, so has world~$1$.
\end{theorem}
\begin{proof}
Take densities. By Theorem~\ref{thm:aexttransfer}, for $m\ge m^*$ world~$1$
is fine on $R_m$ and has an SI pair and an AC pair with non-null parts
inside $R_m$. Proposition~\ref{prop:localext} with $R\coloneqq R_m$ gives a
two-sided extension $P^m_d(\debt\cap R_m)-\kappa_mP^m_c(\cred\cap R_m)$ for
pairs agreeing off $R_m$, with $P^m_d,P^m_c$ positive on non-null sets. In
CBOI's form its charges are $\nu^m_1=\kappa_mP^m_c$ and $\nu^m_2=P^m_d$, up to
a constant, since $\cred=R_m\setminus\{r_1\ge1\}$ and $\debt=\{r_1\ge2\}$ on
$R_m$. Use the sequence $(Z_m)_{m\ge m^*}$.
\end{proof}

\begin{corollary}\label{cor:aedge}
Let $n=3$, let all three rankings be complete with (G) and (IB), assume (S),
and suppose $V$ is undominated. Call $[e_v,e_w]$ an \emph{A-edge} if world $v$,
equivalently world $w$, has (AExt) on $D_{vw}$. At most one edge is an
A-edge. So at least two edges are \emph{leverage edges}: neither of their
worlds has (AExt) there, and a fortiori neither has a two-sided local
extension, a nonconstant local group primary, or the local (FC) there.
\end{corollary}
\begin{proof}
Equivalence is Theorem~\ref{thm:aextshared}. Two distinct edges share a world
$v$, which would have (AExt) on both of its regions; Theorem~\ref{thm:aext}
excludes this. The last clause follows from Lemmas~\ref{lem:FC}
and~\ref{lem:leading}.
\end{proof}

\begin{remark}\label{rem:Q}
In a counterexample, therefore, every world fails (AExt) on at least one of
its edge regions, and on that region its partner fails (AExt) too. This
excludes every non-Archimedean structure that can be confined to sets of
arbitrarily small measure while the rest of the world carries an exchange
rate: heaviness at countably many ultrafilters or cone levels, at a null
compact set of levels, or on finitely many disjoint small sets whose
unions are ordered, over an Archimedean base, like the finite comparative
probability without an agreeing measure of Kraft, Pratt and Seidenberg
\cite{KPS}. The following order-theoretic question, in which
(S) does not appear, would settle P1 if answered negatively: \emph{is there a
complete (G)+(IB) ranking that fails (AExt) on some edge region and has no
nonconstant group primary?} If there is none, every world of a
configuration with $V$ undominated has a nonconstant group primary or (AExt)
on both of its regions, and Theorem~\ref{thm:tame} or
Theorem~\ref{thm:aext} makes $V$ dominated. We know no example.
\end{remark}

\section{Low-cone heaviness}\label{sec:LH}

\begin{definition}\label{def:CN}
At world $v$, a measurable $R\subseteq H'$ is \emph{credit-null} if every
pair $(A,T)$ with $A$ non-null and $T\subseteq R$ is AC, and
\emph{debt-infinitesimal} if every pair $(A,S)$ with $A\subseteq R$ non-null
and $S$ non-null is SI. $\CN_v$ and $\DI_v$ denote the families of such sets;
both are hereditary and invariant under null modifications. We write
$\CN_v\ne0$ if some non-null set is credit-null, and likewise for $\DI_v$.
World $v$ satisfies (LH$_v$) if $\CN_v=0$ and (DH$_v$) if $\DI_v=0$; for an
edge $[e_v,e_w]$, (LH$_{v,w}$) says that no low cone $\Lo vwa$, $a>0$, is
credit-null at~$v$.
\end{definition}

A pair $(\varnothing,T)$ with $T$ non-null is SI by strict (G), so the
non-null debt part in the definition of credit-null is essential.
\emph{Standing for this section:} $n=3$; all three rankings are complete
with (G) and (IB); (S). Nothing is assumed about $V$.

\begin{lemma}\label{lem:excl}
No world has both $\CN_v\ne0$ and $\DI_v\ne0$.
\end{lemma}
\begin{proof}
Let $L\in\CN_v$ and $R\in\DI_v$ be non-null. If $R\setminus L$ is non-null,
put $A\coloneqq R\setminus L$ and $T\coloneqq L$; otherwise $R\subseteq L$ up
to a null set, and we split $R$ into non-null halves $A,T$. The pair $(A,T)$
is SI because $R\in\DI_v$, and AC because $L\in\CN_v$.
\end{proof}

\begin{lemma}\label{lem:conestails}
Let $\{v,w,u\}=\{1,2,3\}$.
\begin{enumerate}[label=(\alph*)]
\item If $\CN_v\ne0$, there is $a>0$ with $\Lo vwa\in\CN_v$ and
$\Lo vua\in\CN_v$.
\item If $\DI_v\ne0$, there is $b<\infty$ with $\Hi vwb\in\DI_v$ and
$\Hi vub\in\DI_v$.
\end{enumerate}
\end{lemma}
\begin{proof}
(a) Let $Q\in\CN_v$ be non-null and split it into non-null disjoint
$Q_1,Q_2$. Lemma~\ref{lem:yard}(b) at world $v$ gives $a_i>0$ such that
every subset of $\Lo vw{a_i}$ and of $\Lo vu{a_i}$ is $\preceq^v_cQ_i$; put
$a\coloneqq\min(a_1,a_2)$, and let $L$ be one of the two cones of level $a$.
Let $(A,T)$ be a pair with $A$ non-null and $T\subseteq L$. One of
$A\setminus Q_1$, $A\setminus Q_2$ is non-null, say $A'\coloneqq A\setminus
Q_i$. Then $T\preceq^v_cQ_i$, and $(A',Q_i)$ is AC since $Q_i\in\CN_v$. In
the context $V+\debt_v(A')$, silent on $Q_i\cup T$,
\[
 V\preceq_vV+\debt_v(A')+\cred_v(Q_i)\preceq_vV+\debt_v(A')+\cred_v(T)
 \preceq_vV+\debt_v(A)+\cred_v(T),
\]
the last step by (G). So $(A,T)$ is AC.

(b) Let $R\in\DI_v$ be non-null and split it into non-null disjoint
$R_1,R_2$. Lemma~\ref{lem:yard}(a) at world $v$ gives $A^i_0\preceq^v_dR_i$
and $b_i$ with $\Hi vw{b_i},\Hi vu{b_i}\prec^v_dA^i_0$; put
$b\coloneqq\max(b_1,b_2)$ and let $H_b$ be one of the two tails of level
$b$. Let $B\subseteq H_b$ and $S$ be non-null and disjoint. If $S\setminus
R_1$ is non-null put $i\coloneqq1$ and $S''\coloneqq S\setminus R_1$;
otherwise $S\subseteq R_1$ up to a null set, so $S$ misses $R_2$: put
$i\coloneqq2$ and $S''\coloneqq S$. Then $S''$ is non-null and misses
$B\cup R_i$, $B\preceq^v_dR_i$, and $(R_i,S'')$ is SI. So
$V+\debt_v(B)+\cred_v(S)\preceq_vV+\debt_v(B)+\cred_v(S'')\preceq_v
V+\debt_v(R_i)+\cred_v(S'')\prec_vV$, by (G) and then (IB) in the context
$V+\cred_v(S'')$. So $(B,S)$ is SI.
\end{proof}

\begin{lemma}\label{lem:EPnotCN}
If (EP$_{vw}$) holds, no low cone $\Lo vwa$ is credit-null at~$v$.
\end{lemma}
\begin{proof}
Fix $0<b<a$ and a non-null $Z\subseteq D_{vw}\cap\{b<t^v_w<a\}$. By CCF's
Lemmas~3 and~4 in world $v$'s frame, with $s^*=b$ approached from below,
there is $q^*$ with $C_{q^*}=E_w(b)$ (debt at $v$ on $\{t^v_w<b\}$, silence
on the rest of $D_{vw}$) such that indices near $q^*$ carry $E_w(s)$ for all
$s$ in some $[s_0,b)$. $E_w(b)$ is silent on $Z$, so
$Y\coloneqq E_w(b)+\cred_v(Z)\prec_vE_w(b)$, and by (T1) $Y\prec_vE_w(s)$ for
some $s\in(0,b)$. Since $Y=E_w(s)+\debt_v(D_{vw}\cap\{s\le t^v_w<b\})
+\cred_v(Z)$, (IB) makes $(D_{vw}\cap\{s\le t^v_w<b\},Z)$ an SI pair with
non-null debt part and credit part inside $\Lo vwa$.
\end{proof}

\begin{lemma}\label{lem:CNtoDI}
Let $E=(\alpha^v_E,\beta^v_E)$ be a selection on the open edge $(e_v,e_w)$
with $\beta^v_E>0$, and let $R\subseteq\cred_v(E)\cap D_{vw}$ be credit-null
at~$v$. Then $R$ is debt-infinitesimal at~$w$.
\end{lemma}
\begin{proof}
Let $A\subseteq R$ and $S$ be non-null and disjoint. Choose a non-null
$S_0\subseteq S$ with $S'\coloneqq S\setminus S_0$ non-null.
Lemma~\ref{lem:yard}(b) at world $w$ gives $a'>0$ such that every subset
of $\Lo wu{a'}$ is $\preceq^w_cS'$. The set
\[
 Q\coloneqq P_w\cap\{t^v_w<\beta^v_E\}\cap\Lo wu{a'}
\]
is non-null (an open cone condition in $P_w$), lies in $\debt_v(E)\cap
D_{vw}$, and misses $A\subseteq\cred_v(E)$. Let $X$ be $E$ silenced on
$Q\cup A$, so $E=X+\debt_v(Q)+\cred_v(A)$. Split $Q$ into non-null $Q_1,Q_2$.
$(Q_1,A)$ is AC because $A\subseteq R\in\CN_v$, so
$V\preceq_vV+\debt_v(Q_1)+\cred_v(A)\prec_vV+\debt_v(Q)+\cred_v(A)$ by strict
(G), and by (IB) $X\prec_vE$. Competitiveness of $E$ on $\{v,w\}$ gives
$E\preceq_wX$. By \eqref{eq:swap}, $E=X+\debt_w(A)+\cred_w(Q)$, so (IB) gives
$V+\debt_w(A)+\cred_w(Q)\preceq_wV$. Since $Q\preceq^w_cS'$ and the context
$V+\debt_w(A)$ is silent on $Q\cup S'$,
\[
 V+\debt_w(A)+\cred_w(S)\prec_wV+\debt_w(A)+\cred_w(S')\preceq_w
 V+\debt_w(A)+\cred_w(Q)\preceq_wV,
\]
the first step by strict (G) on $S_0$. So $(A,S)$ is SI at~$w$.
\end{proof}

\begin{lemma}\label{lem:DItoCN}
Let $v\ne w$. If $\DI_w\ne0$, then $\CN_v\ne0$.
\end{lemma}
\begin{proof}
By Lemma~\ref{lem:conestails}(b), $\Hi wvb\in\DI_w$ for some $b<\infty$. On
the edge $[e_v,e_w]$, in world $v$'s frame, Lemma~\ref{lem:ends}(c) gives a
boundary point $q_0$ of $\{q:C_q=D_v\}$ and $q_m\to q_0$ in the open edge
with $C_m\coloneqq C_{q_m}\ne D_v$ and $\alpha^v_m\to0$. In world $w$'s
frame $C_m$ has thresholds $\alpha^w_m=1/\beta^v_m$ and
$\beta^w_m=1/\alpha^v_m\to\infty$, and $D_v$ is the vertex model that gives
debt at $w$ on all of $D_{vw}$.

\emph{Case I: $\beta^v_m>0$ for infinitely many $m$.} Along them
$\alpha^w_m<\infty$, and Lemma~\ref{lem:SAC} in world $w$'s frame applies at
$q_0$. With $A'\coloneqq\Hi wvb\cap\{t^w_v<b+1\}$, non-null, it gives for
large $m$ a strictly AC pair $(A',\Hi wv{\alpha^w_m})$ at world~$w$. But $A'$
is debt-infinitesimal at $w$, and the tail is non-null and, for
$\alpha^w_m>b+1$, disjoint from $A'$, so the pair is SI: a contradiction.

\emph{Case II: $\beta^v_m=0$ for all large $m$.} Then $C_m=(\alpha^v_m,0)$
with $0<\alpha^v_m<\infty$, silent on $D_{vw}\cap\{t^v_w<\alpha^v_m\}=
\Hi wv{1/\alpha^v_m}$. Fix such $m$ with $1/\alpha^v_m>b$ and put
$T'\coloneqq P_w\cap\{t^v_w<\alpha^v_m\}$, non-null, in the silence of $C_m$,
and in $\Hi wvb$, hence in $\DI_w$. We show $T'\in\CN_v$. Let $T\subseteq T'$
and $A$ be non-null and disjoint. Lemma~\ref{lem:yard}(a) at world $v$ gives
a debt yardstick $S\coloneqq Z_I(A)\subseteq P_v\cap\{t^v_w\in I\}$,
$I\coloneqq(\alpha^v_m/4,\alpha^v_m/2)$, with $S\preceq^v_dA$; $S$ lies in
the silence of $C_m$ and misses $T$ ($P_v$ versus $P_w$). Let
$Y\coloneqq C_m+\debt_w(T)+\cred_w(S)$. Since $T\in\DI_w$, $(T,S)$ is SI at
$w$, so $Y\prec_wC_m$ by (IB), and competitiveness of $C_m$ on $\{v,w\}$ gives
$C_m\preceq_vY$. By \eqref{eq:swap}, $Y=C_m+\cred_v(T)+\debt_v(S)$, so by
(IB) $V\preceq_vV+\debt_v(S)+\cred_v(T)\preceq_vV+\debt_v(A)+\cred_v(T)$, the
second step since $S\preceq^v_dA$ in the context $V+\cred_v(T)$, silent on
$S\cup A$. So $(A,T)$ is AC. A pair with null credit part and non-null debt
part is AC by strict (G). Hence $T'\in\CN_v$.
\end{proof}

\begin{theorem}\label{thm:LH}
Let $n=3$, let all three rankings be complete with (G) and (IB), and assume
(S). Then at every world $v$:
\begin{enumerate}[label=(\roman*)]
\item (LH$_v$): for every non-null $R$ there is an SI pair $(A,T)$ of world
$v$ with $A$ non-null and $T\subseteq R$; in particular (LH$_{v,w}$) holds
for both edges at $v$;
\item (DH$_v$): for every non-null $R$ there is an AC pair $(A,S)$ of world
$v$ with $A\subseteq R$ and $A,S$ non-null.
\end{enumerate}
No assumption on $V$ is used.
\end{theorem}
\begin{proof}
Suppose $\CN_v\ne0$ for some $v$; relabel $v=1$. By
Lemma~\ref{lem:conestails}(a), $\Lo12a$ and $\Lo13a$ are credit-null at
world~$1$ for some $a>0$, so by Lemma~\ref{lem:EPnotCN} neither (EP$_{12}$)
nor (EP$_{13}$) holds: there are two-sided selections $E$ on $(e_1,e_2)$ and
$E'$ on $(e_1,e_3)$. The sets
\[
 R\coloneqq P_1\cap\{t_2>\alpha_E\}\cap\Lo13a,\qquad
 R'\coloneqq P_1\cap\{t_3>\alpha_{E'}\}\cap\Lo12a
\]
are non-null and credit-null at~$1$, and $R\subseteq\cred(E)\cap D_{12}$,
$R'\subseteq\cred(E')\cap D_{13}$, since $P_1\subseteq D_{12}\cap D_{13}$.
Lemma~\ref{lem:CNtoDI} gives $R\in\DI_2$ and $R'\in\DI_3$. Lemma~\ref{lem:DItoCN}
with $(w,v)=(2,3)$ gives $\CN_3\ne0$, which together with $\DI_3\ne0$
contradicts Lemma~\ref{lem:excl}. So $\CN_v=0$ at every world, which is
(LH$_v$). If $\DI_w\ne0$ for some $w$, Lemma~\ref{lem:DItoCN} gives
$\CN_v\ne0$ for $v\ne w$, which has just been excluded.
\end{proof}

The contradiction is reached on the edge $[e_2,e_3]$ opposite the world
where credit-nullity was assumed. On a single edge, credit-nullity at one
world and debt-infinitesimality at the other are mutually consistent:
Lemmas~\ref{lem:CNtoDI} and~\ref{lem:DItoCN} run in opposite directions.
Completeness enters as follows: at world~$1$ through (T1) in
Lemma~\ref{lem:EPnotCN} and in reading AC pairs in
Lemmas~\ref{lem:conestails}(a) and~\ref{lem:CNtoDI}; at the partner worlds
through Lemma~\ref{lem:yard}(a) in Lemma~\ref{lem:conestails}(b) and in
Case~II of Lemma~\ref{lem:DItoCN}, and through Lemma~\ref{lem:SAC} in Case~I.
Lemma~\ref{lem:excl} uses neither (G) nor completeness. Only edge selections
are used.

\begin{corollary}\label{cor:relLH}
Let $\{v,w,u\}=\{1,2,3\}$ and let $E$ be a selection on the open edge
$(e_v,e_w)$.
\begin{enumerate}[label=(\alph*)]
\item If $\beta^v_E>0$, then for every non-null $T\subseteq\cred_v(E)\cap
D_{vw}$ and every $a'>0$ there are non-null $T''\subseteq T$ and
non-null $Q\subseteq P_w\cap\debt_v(E)\cap\Lo wu{a'}$ with $(Q,T'')$ an
SI pair of
world~$v$.
\item If $\beta^v_E<\alpha^v_E$, then for every non-null
$T'\subseteq\sil(E)\cap P_w$ and every $b<\infty$ there is a non-null
$S\subseteq\sil(E)\cap P_v\cap\Hi vub$ with $(T',S)$ an AC pair of
world~$w$.
\end{enumerate}
\end{corollary}
\begin{proof}
(a) Otherwise, for some $T$ and $a'_0$, every pair $(Q,T'')$ with $Q$ a
non-null subset of $Q_0\coloneqq P_w\cap\debt_v(E)\cap\Lo wu{a'_0}$ and
$T''\subseteq T$ is AC at $v$. Run the proof of Lemma~\ref{lem:CNtoDI} with
$R\coloneqq T$, lowering the level $a'$ there to at most $a'_0$ (subsets of
smaller cones are still $\preceq^w_cS'$), so that $Q\subseteq Q_0$; its only
use of credit-nullity is that $(Q_1,A)$ is AC, which is now the assumption.
So $T\in\DI_w$, contradicting Theorem~\ref{thm:LH}.

(b) Otherwise, for some $T'$ and $b$, $(T',S)$ is SI at $w$ for every
non-null $S\subseteq\sil(E)\cap P_v\cap\Hi vub$, and then so is $(T,S)$ for
every non-null $T\subseteq T'$, by (G). Let $T\subseteq T'$ and $A$ be
non-null and disjoint. By Lemma~\ref{lem:yard}(a) at world $v$ there is
$b_A$ such that every subset of $\Hi vu{b_A}$ is $\preceq^v_dA$; let
$S\coloneqq P_v\cap\Hi vu{\max(b,b_A)}\cap\{t^v_w\in I\}$ with $I$ a bounded
open interval inside $(\beta^v_E,\alpha^v_E)$. Then $S$ is non-null, lies in
$\sil(E)$, misses $T$, and $S\preceq^v_dA$. Case~II of the proof of
Lemma~\ref{lem:DItoCN}, with $C_m$ replaced by $E$, gives $(A,T)$ AC at $v$.
So $T'\in\CN_v$, contradicting Theorem~\ref{thm:LH}.
\end{proof}

These are lower Archimedean bounds: credit and debt are never negligible.
They do not give fineness, which is an upper bound; see
Proposition~\ref{ex:germ}.

\begin{corollary}\label{cor:LHcons}
Let $j\in\{2,3\}$ and let $k$ be the third world.
\begin{enumerate}[label=(\alph*)]
\item For every non-null $A$ and all open intervals $I,K\subseteq(0,\infty)$
there is an AC pair $(A'',W)$ of world~$1$ with $A''\preceq_dA$,
$A''\subseteq D_{1j}\cap\{t_j\in I\}$ and $W\subseteq D_{1j}\cap\{t_j\in K\}$,
both non-null.
\item For every non-null $T$ and all open intervals $I,J\subseteq(0,\infty)$
there is an SI pair $(Z,T')$ of world~$1$ with $T'\preceq_cT$,
$T'\subseteq D_{1j}\cap\{t_j\in J\}$ and $Z\subseteq D_{1j}\cap\{t_j\in I\}$,
both non-null.
\item For every non-null $Z\subseteq P_1$ there is $c<\infty$ such that
$V+o_{12}(Y)+o_{23}(Z)\prec_3V$ for every non-null $Y\subseteq P_3\cap
\Hi32c$.
\end{enumerate}
\end{corollary}
\begin{proof}
(a) Let $A'\coloneqq Z_I(A)\subseteq P_1\cap\{t_k>b'\}\cap\{t_j\in I\}$
(Lemma~\ref{lem:yard}(a)). By (DH), $A'$ is not debt-infinitesimal: some pair
$(A'',S)$ with $A''\subseteq A'$ and $S$ non-null is AC. By
Lemma~\ref{lem:yard}(b), $W\coloneqq Z^c_K(S)\subseteq P_1\cap\{t_k<a'\}
\cap\{t_j\in K\}$ satisfies $W\preceq_cS$, and with $a'<b'$ it misses $A''$.
In the context $V+\debt(A'')$, silent on $S\cup W$,
$V\preceq_1V+\debt(A'')+\cred(S)\preceq_1V+\debt(A'')+\cred(W)$. And
$A''\preceq_dA'\preceq_dA$.

(b) Let $T_y\coloneqq Z^c_J(T)\subseteq P_1\cap\{t_k<a'\}\cap\{t_j\in J\}$.
By (LH), some SI pair $(A,T')$ has $A$ non-null and $T'\subseteq T_y$. By
Lemma~\ref{lem:yard}(a), with tail parameter $b'>a'$, $Z\coloneqq Z_I(A)
\subseteq P_1\cap\{t_k>b'\}\cap\{t_j\in I\}$ satisfies $Z\preceq_dA$ and
misses $T'$; in the context $V+\cred(T')$, silent on $Z\cup A$,
$V+\debt(Z)+\cred(T')\preceq_1V+\debt(A)+\cred(T')\prec_1V$. And
$T'\preceq_cT_y\preceq_cT$.

(c) The proof of Lemma~\ref{lem:B1}, with the SI pair of its step~(2) taken
from (LH) at world~$3$ instead of from $\Lambda_3$: by (LH) there is an SI
pair $(A,T'')$ of world~$3$ with $A$ non-null and $T''\subseteq L$, and the
remaining steps are unchanged, with the given set in the role of $B$, the
sets $Y$ in the role of the subsets of $Z_B$, and $c\coloneqq\max(b,a')$. The
debt part $A$ may now meet the given set; this is harmless, since the chain
of step~(4) uses only that $B$ misses $Z_B$ and $T''$.
\end{proof}

\noindent Part (a) supplies, at every level, the small AC pairs whose debt
part is a yardstick and whose credit part sits at prescribed levels; part (b)
supplies the small SI pairs with prescribed levels. These are exactly the
exchange pairs needed to exclude heavy germs at precise points of edge
selections. Part (c) is Lemma~\ref{lem:B1} without a fine edge. In
particular, no set $A$ has the property that $(A',Q)$ is SI for every
$A'\preceq_dA$ and every non-null $Q$: $A$ would be debt-infinitesimal.

\section{Archimedean axioms}\label{sec:arch}

In this section world~$1$ is complete with (G) and (IB); $n\ge2$ unless (S)
is assumed. A regular model is a \emph{cone model} if its credit, silence
and debt sets on $H'$ are invariant under $g\mapsto\lambda g$, $\lambda>0$.
Coherent models, edge models, $V$, $D_i$, the model $W_1$ of
Section~\ref{sec:tame}, with debt at~$1$ on all of $H'$, and their modifications on bands are cone models. Cone
models $M_m$ \emph{converge conically} to $M$ if $\mu\{M_m\ne M\}\to0$; for
edge models this is convergence of the thresholds in $[0,\infty]$.

\begin{definition}\label{def:arch}
World $1$ satisfies
\begin{itemize}[leftmargin=*,itemsep=1pt]
\item (CC-a) if, for cone models $M_m\to M$ conically and every regular $X$,
$M_m\preceq_1X$ for all $m$ implies $M\preceq_1X$; and (CC-b) if
$X\preceq_1M_m$ for all $m$ implies $X\preceq_1M$;
\item (MC) if for every decreasing sequence of cone sets $C_m\subseteq H'$
with null intersection and all non-null $A,T$, some $C_m$ has
$C_m\preceq_dA$ and some has $C_m\preceq_cT$;
\item (GC) (\emph{germ continuity}) if for every non-null $A$ and $T$ and
every germ point of every edge $[e_1,e_j]$, some one-sided band at it is
$\preceq_dA$ and some is $\preceq_cT$;
\item (BI-a), resp.\ (BI-b), if (CC-a), resp.\ (CC-b), holds along the
sequences $V+\cred(b_m)\to V$ and ($W_1$ silenced on $b_m$)$\to W_1$, resp.\
$V+\debt(b_m)\to V$ and ($D_1$ silenced on $b_m$)$\to D_1$, and along their
tails $m\ge m_0$, where $b_m$ is the one-sided band of size $1/m$ at any germ
point of any edge.
\end{itemize}
\end{definition}

(GC) says that the comparative probabilities $\preceq_d,\preceq_c$ are
continuous along every band shrinking to a cone level; it is an Archimedean
axiom in the sense of Savage's P6, restricted to bands. (MC) is monotone
continuity in Villegas' sense \cite{Villegas}, restricted to cones.

\begin{proposition}\label{prop:GC}
\begin{enumerate}[label=(\alph*)]
\item (GC) implies that world~$1$ is fine. No (S) is used.
\item (CC-a)$\Rightarrow$(BI-a)$\Rightarrow$(GC),
(CC-b)$\Rightarrow$(BI-b)$\Rightarrow$(GC), and (MC)$\Rightarrow$(GC).
\end{enumerate}
\end{proposition}
\begin{proof}
(a) Under (GC) no germ point is $A$-heavy or $T$-heavy, so by
Lemma~\ref{lem:germ} every $D_{1j}$ is in $\Jid_A$ and in $\Kid_T$, and so is
$H'=\bigcup_jD_{1j}$.

(b) The sequences in (BI-a), (BI-b) consist of cone models converging
conically, since bands are cones and $\mu(b_m)\to0$; and the bands $b_m$ at a
germ point are decreasing cone sets with null intersection. It remains to
show that (BI-a) and (BI-b) each imply (GC). Suppose (GC) fails: for some
non-null $A$ (debt) and germ point, no band $b_m$ is $\preceq_dA$. Since
$\mu(b_m)\to0$, $A'\coloneqq A\setminus b_{m_0}$ is non-null for some $m_0$.
For $m\ge m_0$, $b_m\cap A'=\varnothing$, and $b_m\not\preceq_dA'$, since
$A'\preceq_dA$; by completeness $A'\prec_db_m$, that is,
$V+\debt(A')\prec_1V+\debt(b_m)$. (BI-b): put $X\coloneqq V+\debt(A')$ and
$M_m\coloneqq V+\debt(b_m)\to V$; then $X\prec_1M_m$ for $m\ge m_0$, so
$X\preceq_1V$, while $V\prec_1X$ by strict (G). (BI-a): let $N_m$ be $W_1$
silenced on $A'\cup b_m$, $M_m\coloneqq N_m+\debt(A')$ ($W_1$ silenced on
$b_m$) and $X\coloneqq N_m+\debt(b_m)$ ($W_1$ silenced on $A'$). By (IB),
$M_m\prec_1X$; $M_m\to W_1$, so $W_1\preceq_1X$, while $X\prec_1W_1$.

Suppose instead that for some non-null $T$ no band $b_m$ is $\preceq_cT$.
With $T'\coloneqq T\setminus b_{m_0}$ non-null, completeness gives
$T'\prec_cb_m$, that is, $V+\cred(b_m)\prec_1V+\cred(T')$, for $m\ge m_0$.
(BI-a): $M_m\coloneqq V+\cred(b_m)\to V$ and $X\coloneqq V+\cred(T')$ give
$V\preceq_1X$, while $X\prec_1V$. (BI-b): let $N_m$ be $D_1$ silenced on
$T'\cup b_m$, $M_m\coloneqq N_m+\cred(T')$ ($D_1$ silenced on $b_m$) and
$X\coloneqq N_m+\cred(b_m)$ ($D_1$ silenced on $T'$). Then $X\prec_1M_m$,
$M_m\to D_1$, so $X\preceq_1D_1$, while $D_1\prec_1X$.

Only strict premises were used, and only targets $X$ that differ from the
limit by a change on a set away from the germ.
\end{proof}

For the endpoint chain of $[e_1,e_j]$ we use CCF's parametrization
$M^j_s\coloneqq E^j_1(-1/s)$ for $s<0$, $M^j_{-\infty}\coloneqq D_1$ and
$M^j_s\coloneqq E_j(s)$ for $s\ge0$; by CCF's equation (2),
$s<s'$ implies $M^j_s\prec_1M^j_{s'}$.

\begin{theorem}\label{thm:ECC}
Let $n\ge3$, assume (G) at every world and (S). Suppose that for every
$j\ne1$ and every regular $X$:
\begin{enumerate}[label=(ECC-\alph*),leftmargin=*]
\item if $s_m\uparrow s$ and $M^j_{s_m}\prec_1X$ for all $m$, then
$M^j_s\preceq_1X$;
\item if $s_m\downarrow s$ and $X\prec_1M^j_{s_m}$ for all $m$, then
$X\preceq_1M^j_s$.
\end{enumerate}
Then world~$1$ is fine, and by Corollary~\ref{cor:ext}(b) it has a loss of
CBOI's form with \eqref{eq:positive} that weakly extends $\preceq_1$ on all
regular pairs. Only targets $X$ that differ from $M^j_s$ by a yardstick
change at levels away from the limit level are used.
\end{theorem}
\begin{proof}
Suppose a germ point of some edge $[e_1,e_j]$ is $A$-heavy or $T$-heavy. By
Corollary~\ref{cor:X5} it is not $(\infty,-)$ for debt or $(0,+)$ for credit.
If $Z\preceq_dA$, heaviness and Lemma~\ref{lem:ideal} give $b_m\notin\Jid_Z$,
so $Z\prec_db_m$ by completeness, for every one-sided band $b_m$ of the germ
point; likewise $Z\prec_cc_m$ if $Z\preceq_cT$. Bands $b_m,c_m$ below are of
size $1/m$.

\emph{Debt, $(z,+)$, $0\le z<\infty$.} $M^j_{z+1/m}=E_j(z)+\debt(b_m)$ with
$b_m=\Band_j[z,z+1/m)$. Let $Z\coloneqq Z_I(A)$ with $I\coloneqq(z+1,z+2)$,
silent for all these models, and $X\coloneqq E_j(z)+\debt(Z)$. The pairs
$(X,M^j_{z+1/m})$ and $(V+\debt(Z),V+\debt(b_m))$ differ exactly on
$Z\cup b_m$ with the same advice, so $X\prec_1M^j_{z+1/m}$; (ECC-b) gives
$X\preceq_1E_j(z)$, but $E_j(z)\prec_1X$ by strict (G).

\emph{Debt, $(z,-)$, $0<z<\infty$.} For large $m$,
$M^j_{z-1/m}$ is $E_j(z)$ silenced on $b_m=\Band_j[z-1/m,z)$. Let
$Z\coloneqq Z_I(A)$ with $I\coloneqq(z/4,z/2)$, in the debt of all these
models for large $m$, and let $X$ be $E_j(z)$ silenced on $Z$. Then
$M^j_{z-1/m}\prec_1X$ by (IB), (ECC-a) gives $E_j(z)\preceq_1X$, but
$X\prec_1E_j(z)$.

\emph{Credit, $(w,-)$, $0<w\le\infty$.} $E^j_1(w-1/m)=E^j_1(w)+\cred(c_m)$
with $c_m=\Band_j(w-1/m,w]$ (for $w=\infty$, $E^j_1(m)=E^j_1(\infty)+
\cred(\Band_j(m,\infty))$); the parameters $-1/(w-1/m)$ increase to $-1/w$
(to $0$ if $w=\infty$). Let $Z\coloneqq Z^c_I(T)$ with $I\coloneqq(w/4,w/2)$
($I\coloneqq(1,2)$ if $w=\infty$), silent for these models for large $m$, and
$X\coloneqq E^j_1(w)+\cred(Z)$. Then $E^j_1(w-1/m)\prec_1X$ by (IB) and
$Z\prec_cc_m$; (ECC-a) gives $E^j_1(w)\preceq_1X$, but $X\prec_1E^j_1(w)$.

\emph{Credit, $(w,+)$, $0<w<\infty$.} $E^j_1(w+1/m)$ is $E^j_1(w)$ silenced
on $c_m=\Band_j(w,w+1/m]$, and the parameters decrease to $-1/w$. Let
$Z\coloneqq Z^c_I(T)$ with $I\coloneqq(w+1,w+2)$, in the credit of all these
models, and let $X$ be $E^j_1(w)$ silenced on $Z$. Then
$X\prec_1E^j_1(w+1/m)$; (ECC-b) gives $X\preceq_1E^j_1(w)$, but
$E^j_1(w)\prec_1X$.

So no germ point is heavy, and Lemma~\ref{lem:germ} gives fineness.
\end{proof}

Within the endpoint chain a band can vanish from the worse side only as a
right band at the limit threshold and from the better side only as a left
band, so (ECC-a) handles exactly the left germs and (ECC-b) exactly the right
germs. Under (EP$_{1j}$), CCF's Lemma~3 makes $s\mapsto M^j_s$ the image of a
continuous surjection $q\mapsto s(q)$ of the edge with $C_q=M^j_{s(q)}$; then
$\{s:M^j_s\preceq_1X\}=s(A_1(X)\cap[e_1,e_j])$ is compact, which is (ECC-a)
for all $X$.

\begin{theorem}\label{thm:arch}
Let $n=3$, let all three rankings be complete with (G) and (IB), and assume
(S). If one world satisfies (GC) or any of (CC-a), (CC-b), (BI-a),
(BI-b), (MC), which imply it, or (ECC-a) and (ECC-b) on each of its edges, or (EP) on each
of its edges, then $V$ is dominated. For any $n\ge3$, if all rankings are
complete with (G) and (IB), (S) holds, and every world satisfies one of these
conditions, then every coherent model whose credal set contains three
vertices is dominated.
\end{theorem}
\begin{proof}
Each condition makes the world fine: Proposition~\ref{prop:GC},
Theorem~\ref{thm:ECC}, CCF's Theorem~7. For $n=3$ apply
Theorem~\ref{thm:onefine}. For general $n$, Corollary~\ref{cor:ext}(b) at
every world and CCF's Lemma~1.
\end{proof}

The axioms are not consequences of completeness, (G) and (IB), even for a
world that is ranked exactly by a loss of CBOI's form.

\begin{proposition}\label{ex:germ}
Fix $z\in(0,\infty)$. Let $\Uf_0$ be an ultrafilter containing the bands
$\Band_2(z,z+1/m)$, $m\ge1$ (they decrease and are non-null), and put
$\pi(B)\coloneqq1$ if $B\cap H'\in\Uf_0$ and $\pi(B)\coloneqq0$ otherwise.
Let world~$1$ be ranked by the loss of CBOI's form with
$\nu_{11}=\nu_{12}\coloneqq\mu+\pi$ and $\nu_{13}\coloneqq\mu$. Then:
\begin{enumerate}[label=(\alph*)]
\item world~$1$'s ranking is complete and satisfies (G) and (IB), and its
loss satisfies \eqref{eq:positive} and weakly, indeed exactly, extends it;
\item (GC) fails at $(z,+)$ on $[e_1,e_2]$, and world~$1$ is not fine on
$D_{12}$; so (CC-a), (CC-b), (MC) fail;
\item (LH$_1$) and (DH$_1$) hold.
\end{enumerate}
\end{proposition}
\begin{proof}
$\pi$ is a $\{0,1\}$-valued finitely additive charge vanishing on null sets:
of two disjoint sets at most one is in $\Uf_0$, and if their union is in
$\Uf_0$, exactly one is. Put $\nu\coloneqq\mu+\pi$ on $H'$.

(a) The loss is real-valued, so the ranking is complete. If $M'$ is a weak
pointwise improvement of $M$, every rank set $\{r_1\ge l\}$ shrinks, so
$L(M')\le L(M)$; if the improvement is strict on a non-null $Q$, some
$Q_l\coloneqq Q\cap\{r_1(\cdot,M')<l\le r_1(\cdot,M)\}$ is non-null and
$L(M)-L(M')\ge\mu(Q_l)>0$. This is (G). By finite additivity
$L(M)-L(M')$ depends only on the advice of $M$ and $M'$ where they differ,
which is (IB).

(b) For $B,C\subseteq H'$, $L(V+\debt(B))-L(V+\debt(C))=\nu(B)-\nu(C)$, so
$B\preceq_dC$ iff $\nu(B)\le\nu(C)$. Let $A\coloneqq\Band_2(z+1,z+2)$. It is
disjoint from $\Band_2(z,z+1)\in\Uf_0$, so $\pi(A)=0$ and $\nu(A)=\mu(A)<1$,
while $\nu(b_m)\ge1$ for every band $b_m$ at $(z,+)$. So no $b_m$ is
$\preceq_dA$. If $D_{12}$ were contained in a finite union of sets
$\preceq_dA$, one of them
would be in $\Uf_0$ and have $\nu\ge1>\nu(A)$. The other axioms imply (GC)
by Proposition~\ref{prop:GC}.

(c) A pair $(A,T)$ is SI iff $\nu(A)<\nu(T)$. Given a non-null $R$, split it
into non-null $R_1,R_2$; one of them, say $R_1$, is not in $\Uf_0$. Take
$T\coloneqq R_2$ and a non-null $A\subseteq R_1$ with $\mu(A)<\mu(R_2)$; then
$\pi(A)=0$ and $\nu(A)<\nu(T)$, which proves (LH$_1$). For (DH$_1$) take a
non-null $A\subseteq R$ whose complement $H'\setminus A$ is non-null, split the
complement into two non-null parts, one of them not in $\Uf_0$, and take in it
a non-null $S$ with $\mu(S)\le\mu(A)$; then $\nu(S)\le\nu(A)$ and $(A,S)$ is
AC.
\end{proof}

So violating the Archimedean axioms is not the same as carrying leverage:
this world has an exact finitely additive extension. By
Theorem~\ref{thm:tame}(ii), if world~$1$ of a three-world configuration with
complete (G)+(IB) rankings and (S) is of this kind, $V$ is dominated; and by
Corollary~\ref{cor:edgefine}, world~$2$ cannot be fine on $D_{12}$ in such a
configuration. Whether a world of this kind is compatible with (S) at all is
open. Part (c) shows that (LH) and (DH) do not imply fineness.

\section{The shape of a counterexample}\label{sec:shape}

\emph{Standing for this section:} $n=3$; all three rankings are complete
with (G) and (IB); (S); and $V$ is undominated. We use the notation
$T_0(C)$, $B_0(C)$ of Section~\ref{sec:tame} and $\tau_{T,B}$ of
Section~\ref{sec:edge}.

\begin{theorem}\label{thm:ownregion}
Let $E$ be a two-sided selection on $(e_1,e_2)$ and $T\subseteq T_0(E)$
non-null. Then $T$ fails
\begin{enumerate}[label=(M2$^*$),leftmargin=*]
\item for every $c<\infty$, $T$ is a finite union of sets $Q$ such that
$(Q,P_3\cap\Hi32c)$ is a weakly SI pair of world~$2$.
\end{enumerate}
\end{theorem}
\begin{proof}
On $P_2$, $o_{13}$ is debt at $2$, and on $P_3$, $o_{12}$ is credit at $2$;
so (M2$^*$) says that world~$2$'s debt on small pieces of $T$ is compensated by
its credit on deep tails of $P_3$. Suppose (M2$^*$) holds.

\emph{Step 1: for every non-null $Z\subseteq B_0(E)$,
$V\prec_1V+\tau_{T,Z}$.} By Corollary~\ref{cor:LHcons}(c) there is $c$ such
that $V+o_{12}(Y)+o_{23}(Z)\prec_3V$ for $Y\coloneqq P_3\cap\Hi32c$. By
(M2$^*$), $T=Q_1\cup\dots\cup Q_r$ with $(Q_k,Y)$ weakly SI at world~$2$;
replacing $Q_k$ by $Q_k\setminus(Q_1\cup\dots\cup Q_{k-1})$ keeps this, by
(G). Let $T_k\coloneqq Q_1\cup\dots\cup Q_k$ and suppose
$V+\tau_{T,Z}\preceq_1V$. As in the proof of Lemma~\ref{lem:B2}, there is a
least $k^*\ge1$ with $V+\tau_{T_{k^*},Z}\preceq_1V$, and
$V+\tau_{T_{k^*-1},Z}\preceq_2V$ by Lemma~\ref{lem:XOR}. Let
$d\coloneqq o_{13}(T_{k^*})+o_{23}(Z)+o_{12}(Y)$. Then
$V+d\prec_1V+\tau_{T_{k^*},Z}\preceq_1V$ by strict (G) (credit at $1$ on $Y$);
$V+d\preceq_2V+\tau_{T_{k^*-1},Z}\preceq_2V$ by (IB), since the first pair
differs from $(V+\debt_2(Q_{k^*})+\cred_2(Y),V)$ only in shared advice; and
$V+d\preceq_3V+o_{12}(Y)+o_{23}(Z)\prec_3V$ by (G). So $V+d$ dominates $V$, a
contradiction.

\emph{Step 2.} $T$ is credit-null at world~$1$: let $A$ be non-null and
$T''\subseteq T$ disjoint from $A$, and suppose $(A,T'')$ is SI. By
Lemma~\ref{lem:yard}(a) on $[e_1,e_3]$, $Z\coloneqq B_0(E)\cap
\Band_3(b,\infty)$ is non-null and $\prec_dA$ for suitable $b$; in the context
$V+\cred(T'')$, silent on $Z\cup A$, $V+\debt(Z)+\cred(T'')\prec_1
V+\debt(A)+\cred(T'')\prec_1V$, and adding credit on $T\setminus T''$, which
misses $Z$, gives $V+\tau_{T,Z}\prec_1V$, contradicting Step~1. So $T$ is a
non-null credit-null set, contrary to Theorem~\ref{thm:LH}.
\end{proof}

Exchanging the labels $1$ and $2$ maps $T_0(E)$ to $B_0(E)$ and gives the
same statement for world~$1$'s debt on subsets of $B_0(E)$ relative to
$P_3\cap\Hi31c$; relabelling gives it on every edge. So at every two-sided
selection of every edge, each edge world's debt order is, in the sense of
(M2$^*$), non-fine on its own odd region relative to credit on deep tails of
the third region. The proof
is that of Lemmas~\ref{lem:B2} and~\ref{lem:B3}; the fine edge of
Section~\ref{sec:edge} supplied (M2$^*$) there, and Theorem~\ref{thm:LH}
supplies the third world's acceptance (Corollary~\ref{cor:LHcons}(c)) and
excludes the credit-null alternative.

\begin{theorem}\label{thm:dense}
Let $C$ be a two-sided selection on $(e_u,e_v)$, $T_0\coloneqq T_0(C)$ and
$B_0\coloneqq B_0(C)$. For every non-null $Q\subseteq B_0$ there are
non-null $T\subseteq T_0$ and $S\subseteq P_w$ such that $Q\notin\Kid^v_{S'}$
for every $S'\prec^v_cS$. Hence some ultrafilter contains $Q$ and avoids
every $\Kid^v_{S'}$ with $S'\prec^v_cS$: world $v$ is credit-heavy inside
every non-null subset of $B_0$. Symmetrically, world $u$ is credit-heavy in
the same sense inside every non-null subset of $T_0$.
\end{theorem}
\begin{proof}
On $P_u$, $o_{vw}$ is credit at $v$; on $P_v$, $o_{uw}$ is debt at $v$ and
credit at $w$; on $P_w$, $o_{uv}$ is debt at $w$ and credit at $v$.

(1) By (LH$_v$), some SI pair $(A,Q_1)$ of world $v$ has $A$ non-null and
$Q_1\subseteq Q$. In world $v$'s frame $T_0=P_v\cap\{t^v_u<1/\alpha^u_C\}$ meets
every high tail $\Hi vwb$ in a non-null set; by Lemma~\ref{lem:yard}(a) at
world $v$ choose $b$ with every subset of $\Hi vwb$ $\preceq^v_dA$, and put
$T\coloneqq T_0\cap\Hi vwb$. Then $(T,Q_1)$ is SI (context
$V+\cred_v(Q_1)$, silent on $T\cup A$), and so is $(T,Q)$ by (G):
$V+o_{uw}(T)+o_{vw}(Q)\prec_vV$.

(2) By (LH$_w$), some SI pair $(A',T_1)$ of world $w$ has $A'$ non-null and
$T_1\subseteq T$. By Lemma~\ref{lem:yard}(a) at world $w$ choose a non-null
$S\subseteq P_w$ inside a high tail of an edge at $w$ with $S\preceq^w_dA'$.
Then $(S,T_1)$ is SI at $w$, in the context $V+\cred_w(T_1)$, silent on
$S\cup A'$, as in (1), and so is $(S,T)$ by (G):
$V+o_{uv}(S)+o_{uw}(T)\prec_wV$.

(3) Suppose $Q=Q^1\cup\dots\cup Q^r$ with $Q^i\preceq^v_cS'$ and
$S'\prec^v_cS$; we may take the $Q^i$ disjoint. Choose an increasing family
$\Sigma_s\subseteq Q$, $s\in[0,1]$, with $\Sigma_0=\varnothing$,
$\Sigma_1=Q$, which for $s\in[(k-1)/r,k/r]$ fills $Q^k$ continuously in
measure after $Q^1,\dots,Q^{k-1}$; then for every $s^*$ and every $s$ close
to $s^*$ the set $\Sigma_s\setminus\Sigma_{s^*}$ (if $s>s^*$) or
$\Sigma_{s^*}\setminus\Sigma_s$ (if $s<s^*$) lies in a single $Q^i$. Put
$X'\coloneqq V+o_{uv}(S)$, $X_s\coloneqq X'+o_{uw}(T)+o_{vw}(\Sigma_s)$,
$U\coloneqq\{s:X_s\preceq_vX'\}$ and $D\coloneqq\{s:V\preceq_vX_s\}$. Since
$X_s$ improves at $v$ as $s$ grows, $U$ is an up-set and $D$ a down-set.
$U\cap D=\varnothing$, since $X'\prec_vV$ by strict (G). $0\notin U$, since
$X_0$ adds debt at $v$ on $T$ to $X'$. $1\notin D$, since
$X_1\preceq_vV+o_{uw}(T)+o_{vw}(Q)\prec_vV$ by (G) and (1). If some $s$ lies
in neither, then $V\not\preceq_vX_s$, $X_s\not\preceq_vX'$, and
$X_s\preceq_wV+o_{uv}(S)+o_{uw}(T)\prec_wV$ by (G) and (2), and
Lemma~\ref{lem:chain} with $A\coloneqq T\cup\Sigma_s$ (on which $C$ gives
$o_{uw}$, resp.\ $o_{vw}$) gives a contradiction. Otherwise $U\cup D=[0,1]$
with a cut $s^*=\sup D=\inf U$. If $s^*\in D$, take $s>s^*$ close to $s^*$
and $\Delta\coloneqq\Sigma_s\setminus\Sigma_{s^*}$: then
$X'+o_{uw}(T)+o_{vw}(\Sigma_{s^*})+\cred_v(\Delta)\preceq_vX'$ and
$V\preceq_vX_{s^*}$; removing the shared advice $o_{uv}$ on $S$ by (IB), and
writing $\tau\coloneqq o_{uw}(T)+o_{vw}(\Sigma_{s^*})$,
\[
 V+\tau+\cred_v(\Delta)\preceq_vV\preceq_vV+\tau+\cred_v(S),
\]
so $S\preceq^v_c\Delta$ by (IB) in the context $V+\tau$. If $s^*\in U$, the
same computation with $s<s^*$ and $\Delta\coloneqq\Sigma_{s^*}\setminus
\Sigma_s$ gives $S\preceq^v_c\Delta$. Either way
$S\preceq^v_c\Delta\preceq^v_cQ^i\preceq^v_cS'\prec^v_cS$, a contradiction.

The ideals $\Kid^v_{S'}$, $S'\prec^v_cS$, are totally ordered by inclusion,
by completeness, so their union is an ideal, which does not contain $Q$; the
argument of Lemma~\ref{lem:uf}(b) gives the ultrafilter. The symmetric
statement is the same with $u$ and $v$ exchanged, which maps $B_0$ to $T_0$.
\end{proof}

\begin{theorem}\label{thm:shape}
Let $n=3$, let all three rankings be complete with (G) and (IB), assume (S),
and suppose $V$ is undominated. Then:
\begin{enumerate}[label=(\alph*),leftmargin=*,itemsep=2pt]
\item No world is fine (Theorem~\ref{thm:onefine}). Every world carries a
leverage sequence, has no nonconstant group primary, and violates some
finite cancellation condition (Corollary~\ref{cor:tameshape}). Every world
fails (AExt) on at least one of its edge regions, and at most one edge is an
A-edge (Corollary~\ref{cor:aedge}).
\item No edge is fine, and every edge carries a two-sided selection
(Corollary~\ref{cor:fineedge}). Where a world has a two-sided local
extension on an edge region, every heavy ultrafilter of its partner there is
an atom of its charges (Corollary~\ref{cor:tameshape}).
\item Every world satisfies (LH) and (DH) (Theorem~\ref{thm:LH}), with the
relative forms of Corollary~\ref{cor:relLH} and the exchange pairs of
Corollary~\ref{cor:LHcons}.
\item At every two-sided selection $C$ on every edge $(e_u,e_v)$, world $v$
is credit-heavy inside every non-null subset of $B_0(C)$ and world $u$ inside
every non-null subset of $T_0(C)$ (Theorem~\ref{thm:dense}); and each of the
two edge worlds has a debt order that is non-fine on its own odd region,
relative to credit on deep tails of the third region
(Theorem~\ref{thm:ownregion}).
\end{enumerate}
\end{theorem}
\begin{proof}
Collected from the cited results.
\end{proof}

\subsection*{What remains open}

\begin{enumerate}[leftmargin=*,itemsep=3pt]
\item Problem~P1 remains open: we neither prove that $V$ is dominated nor
know a configuration satisfying all conditions of Theorem~\ref{thm:shape}.
\item By Theorem~\ref{thm:ownregion}, P1 follows if (S) forces (M2$^*$) for
one non-null $T\subseteq T_0(E)$ at one two-sided selection $E$ of one edge;
and if some edge has no two-sided selection, it has (EP), and $V$ is
dominated by Lemma~\ref{lem:EPfine} and Corollary~\ref{cor:fineedge}. Every
transfer used in this note moves the debt side of a comparison into the
partner's odd region (Lemmas~\ref{lem:CNtoDI} and~\ref{lem:DItoCN}), so
(M2$^*$), which concerns a world's own odd region, needs a new transfer.
\item The order-theoretic question of Remark~\ref{rem:Q}: a complete
(G)+(IB) ranking that fails (AExt) on an edge region and has no nonconstant
group primary. A negative answer settles P1.
\item For $n\ge4$, Corollary~\ref{cor:edgefine} transfers fineness along
every edge, but the analogue of Proposition~\ref{prop:third}, and with it the
analogue of Theorem~\ref{thm:onefine}, has not been checked. The same holds
for Theorem~\ref{thm:LH}, whose proof uses $n=3$ only through the regions
$P_x$ and the yardsticks.
\item Completeness is used at every world: through (T1), the Savage--Wakker
step, and the reading of AC pairs. For incomplete rankings the results of
this note are not available.
\item Whether a world as in Proposition~\ref{ex:germ}, with a heavy germ
charge, is compatible with (S) at all.
\end{enumerate}

\begin{thebibliography}{9}
\bibitem{Holder} O.~H\"older, Die Axiome der Quantit\"at und die Lehre vom
Mass, \emph{Berichte \"uber die Verhandlungen der K\"oniglich
S\"achsischen Gesellschaft der Wissenschaften zu Leipzig,
Mathematisch-Physische Classe} 53 (1901), 1--64.
\bibitem{KPS} C.~H.~Kraft, J.~W.~Pratt and A.~Seidenberg, Intuitive
probability on finite sets, \emph{Annals of Mathematical Statistics} 30
(1959), 408--419.
\bibitem{Motzkin} T.~S.~Motzkin, \emph{Beitr\"age zur Theorie der linearen
Ungleichungen}, Dissertation, Basel, 1936.
\bibitem{PP} T.~Pennanen and A.-P.~Perkki\"o, Topological duals of locally
convex function spaces, \emph{Positivity} 26, article 2 (2022), Theorem 1.
\bibitem{Savage} L.~J.~Savage, \emph{The Foundations of Statistics}, Wiley,
New York, 1954. Chapter~3, \S3.
\bibitem{Villegas} C.~Villegas, On qualitative probability
$\sigma$-algebras, \emph{Annals of Mathematical Statistics} 35 (1964),
1787--1796.
\bibitem{Wakker} P.~Wakker, Agreeing probability measures for comparative
probability structures, \emph{Annals of Statistics} 9 (1981), 658--662.
\bibitem{YH} K.~Yosida and E.~Hewitt, Finitely additive measures,
\emph{Transactions of the American Mathematical Society} 72 (1952), 46--66.
\end{thebibliography}

\end{document}
